 \documentclass[12pt]{amsart}

\newcommand\hmmax{0}
\newcommand\bmmax{0}

\pdfoutput=1
\usepackage{latexsym, amsmath, amssymb, amsthm, mathrsfs, bm, mathtools,comment, relsize}

\usepackage[table,xcdraw]{xcolor}
\usepackage{subfigure}

\usepackage[mathscr]{eucal}
%\usepackage{showkeys}
%\usepackage{paralist}
\usepackage{array}
\usepackage{tikz}
\usetikzlibrary{arrows,shapes,positioning,calc,patterns,cd}
\usepackage{aliascnt}
%\usetikzlibrary{decorations.pathreplacing}
\usepackage{graphicx} 
%\usepackage[alphabetic]{amsrefs}
\usepackage[T1]{fontenc} 
\usepackage{mathptmx}
\usepackage{microtype}
\renewcommand{\baselinestretch}{1.15}
\usepackage[colorlinks=true,linkcolor=blue,urlcolor=blue, citecolor=blue]{hyperref}
  
  \usepackage[inner=2.3cm,outer=2.3cm, bottom=3.3cm]{geometry}

%\newcommand\hmmax{0}
%\newcommand\bmmax{0}

%\usepackage{mathptmx, amscd, amssymb, amsmath, enumerate, slashbox, %hyperref,float, mathtools}
%\def\baselinestretch{1.15}
%\usepackage{eulervm}
%\usepackage{paralist}
%\usepackage{aliascnt}
%\usepackage{blkarray}
%\usepackage{mathbbol}
%\usepackage[inner=2.3cm,outer=2.3cm, bottom=3.3cm]{geometry}
%\usepackage{setspace}
%\usepackage[colorlinks=true,linkcolor=blue,urlcolor=blue,pagebackref]{hyperref}
\usepackage[initials]{amsrefs}
%\usepackage{marginnote}
%\usepackage{tikz}
%\usetikzlibrary{matrix}
%\usetikzlibrary{arrows,calc}
\allowdisplaybreaks


%\setcitestyle{numbers}
\usepackage{xcolor, color, soul}
\newcommand{\mathcolorbox}[2]{\colorbox{#1}{$\displaystyle #2$}}
%\usepackage[colorlinks=true,pagebackref,hyperindex,citecolor=green,linkcolor=red]{hyperref}

\usepackage{color}
\usepackage[all]{xy}  
\usepackage{enumerate}
\usepackage{mathbbol,mathtools}

%\usepackage{fullpage}
%\usepackage[top=1in, bottom=1in, left=1in, right=1in]{geometry}
%\usepackage{mathrsfs}
\newcommand{\cb}{\color {blue}}                    
\newcommand{\crd}{\color {red}}                   
\newcommand{\cg}{\color {green}}  
%\usepackage{stmaryrd}
\renewcommand{\leq}{\leqslant}
\renewcommand{\geq}{\geqslant}

%\usepackage{mathabx}

% THEOREM Environments ---------------------------------------------------

%%************* New Theorems***************************************
%% Italic text theorems 
\newtheorem{theorem}{Theorem}[section]
\newtheorem{data}{Data}
\providecommand*{\dataautorefname}{Data}
\renewcommand*{\thedata}{\Roman{data}}
\newtheorem{headthm}{Theorem}
\providecommand*{\headthmautorefname}{Theorem}
\renewcommand*{\theheadthm}{\Alph{headthm}}
\newaliascnt{headcor}{headthm}
\newtheorem{headcor}[headcor]{Corollary}
\aliascntresetthe{headcor}
\providecommand*{\headcorautorefname}{Corollary}
\newaliascnt{headconj}{headthm}
\newtheorem{headconj}[headconj]{Conjecture}
\aliascntresetthe{headconj}
\providecommand*{\headconjautorefname}{Conjecture}
\newaliascnt{corollary}{theorem}
\newtheorem{corollary}[corollary]{Corollary}
\aliascntresetthe{corollary}
\providecommand*{\corollaryautorefname}{Corollary}

\newaliascnt{claim}{theorem}
\newtheorem{claim}[claim]{Claim}
\aliascntresetthe{claim}
\providecommand*{\claimautorefname}{Claim}

\newaliascnt{lemma}{theorem}
\newtheorem{lemma}[lemma]{Lemma}
\aliascntresetthe{lemma}
\providecommand*{\lemmaautorefname}{Lemma}


\newaliascnt{conjecture}{theorem}
\newtheorem{conjecture}[conjecture]{Conjecture}
\aliascntresetthe{conjecture}
\providecommand*{\conjectureautorefname}{Conjecture}
\newaliascnt{proposition}{theorem}
\newtheorem{proposition}[proposition]{Proposition}
\aliascntresetthe{proposition}
\providecommand*{\propositionautorefname}{Proposition}



%% Non italic text theorems 
\theoremstyle{definition}
\newaliascnt{definition}{theorem}
\newtheorem{definition}[definition]{Definition}
\aliascntresetthe{definition}
\providecommand*{\definitionautorefname}{Definition}
\newaliascnt{notation}{theorem}
\newtheorem{notation}[notation]{Notation}
\aliascntresetthe{notation}
\providecommand*{\notationautorefname}{Notation}
\newaliascnt{example}{theorem}
\newtheorem{example}[example]{Example}
\aliascntresetthe{example}
\providecommand*{\exampleautorefname}{Example}
\newaliascnt{examples}{theorem}
\newtheorem{examples}[examples]{Examples}
\aliascntresetthe{examples}
\providecommand*{\examplesautorefname}{Examples}
\newaliascnt{remark}{theorem}
\newtheorem{remark}[remark]{Remark}
\aliascntresetthe{remark}
\providecommand*{\remarkautorefname}{Remark}
\newaliascnt{fact}{theorem}
\newtheorem{fact}[fact]{Fact}
\aliascntresetthe{fact}
\providecommand*{\factautorefname}{Fact}
\newaliascnt{question}{theorem}
\newtheorem{question}[question]{Question}
\aliascntresetthe{question}
\providecommand*{\questionautorefname}{Question}
\newaliascnt{questions}{theorem}
\newtheorem{questions}[questions]{Questions}
\aliascntresetthe{questions}
\providecommand*{\questionsautorefname}{Questions}
\newaliascnt{problem}{theorem}
\newtheorem{problem}[problem]{Problem}
\aliascntresetthe{problem}
\providecommand*{\problemautorefname}{Problem}
\newaliascnt{construction}{theorem}
\newtheorem{construction}[construction]{Construction}
\aliascntresetthe{construction}
\providecommand*{\constructionautorefname}{Construction}
\newaliascnt{setup}{theorem}
\newtheorem{setup}[setup]{Setup}
\aliascntresetthe{setup}
\providecommand*{\setupautorefname}{Setup}
\newaliascnt{algorithm}{theorem}
\newtheorem{algorithm}[algorithm]{Algorithm}
\aliascntresetthe{algorithm}
\providecommand*{\algorithmautorefname}{Algorithm}
\newaliascnt{observation}{theorem}
\newtheorem{observation}[observation]{Observation}
\aliascntresetthe{observation}
\providecommand*{\observationautorefname}{Observation}
\newaliascnt{defprop}{theorem}
\newtheorem{defprop}[defprop]{Definition-Proposition}
\aliascntresetthe{defprop}
\providecommand*{\defpropautorefname}{Definition-Proposition}
\newtheorem*{acknowledgement}{Acknowledgement}




\def\equationautorefname~#1\null{(#1)\null}
\def\sectionautorefname~#1\null{Section #1\null}
\def\subsectionautorefname~#1\null{\S #1\null}
\def\Itemautorefname{Part}
\def\trdeg{{\rm trdeg}}
\def\surjects{\twoheadrightarrow}
\def \Fl{\mathcal{F}\ell}
\def \GL{{\operatorname{GL}}}
%\renewcommand*{\qedsymbol}{\(\blacksquare\)}

\newcommand{\luis}[1]{{\color{purple} \sf $\clubsuit$ Luis: [#1]}}
\newcommand{\jonathan}[1]{{\color{purple} \sf $\clubsuit$ Jonathan: [#1]}}
\newcommand{\alessandro}[1]{{\color{purple} \sf $\clubsuit$ Alessandro:[#1]}}
\newcommand{\lisa}[1]{{\color{purple} \sf $\clubsuit$ Lisa: [#1]}}
\newcommand{\matteo}[1]{{\color{purple} \sf $\clubsuit$ Matteo: [#1]}}

% MATH -------------------------------------------------------------------

%For p^e-th roots
\newcommand{\R}[1]{R^{1/p^{#1}}}
\newcommand{\m}{\mathfrak{m}}
\newcommand{\n}{\mathfrak{n}}
\newcommand{\fpt}{{\operatorname{fpt}}}
\newcommand{\grfpt}{{\operatorname{fpt}}_{\\gr}}
\newcommand{\fit}[1]{{\operatorname{fit}}(#1)}

%For sets
\newcommand{\set}[1]{\left\{ #1 \right\}}
\renewcommand{\(}{\left(}
\renewcommand{\)}{\right)}
\newcommand{\RR}{\mathbb{R}}
\newcommand{\NN}{\mathbb{N} }
\newcommand{\ZZ}{\mathbb{Z}}
\newcommand{\QQ}{\mathbb{Q}}
\newcommand{\FF}{\mathbb{F}}
\newcommand{\CC}{\mathbb{C}}

\newcommand{\kk}{\mathbb{k}}

\newcommand{\fn}{\mathbf{n}}
\newcommand{\fx}{\mathbf{x}}
\newcommand{\fa}{\mathbf{a}}
\newcommand{\fb}{\mathbf{b}}

\newcommand{\J}{\mathcal{J}}
\newcommand{\IN}{\operatorname{in}}	

\newcommand{\cK}{\mathcal{K}}
\newcommand{\cR}{\mathcal{R}}
\newcommand{\cV}{\mathcal{V}}
\newcommand{\fp}{\mathfrak{p}}

\newcommand{\egHK}{\ee_{\gHK}}
\newcommand{\cF}{\mathcal{F}}
\newcommand{\cM}{\mathcal{M}}
\newcommand{\cP}{\mathcal{P}}
\newcommand{\Cech}{ \check{\rm{C}}}

\DeclareMathOperator{\leadmon}{{lm}}
\DeclareMathOperator{\mcm}{{mcm}}
\newcommand{\ee}{\operatorname{e}}
\newcommand{\reg}{\operatorname{reg}}
\newcommand{\pd}{\operatorname{pd}}
\newcommand{\sdim}{\operatorname{sdim}}
\newcommand{\Rank}{\operatorname{rank}}
\newcommand{\type}{\operatorname{type}}
\newcommand{\syz}{\operatorname{syz}}
\newcommand{\gHK}{\operatorname{gHK}}
\newcommand{\Spec}{\operatorname{Spec}}
\newcommand{\Depth}{\operatorname{depth}}
\newcommand{\Hom}{\operatorname{Hom}}
\newcommand{\End}{\operatorname{End}}
\newcommand{\gr}{\operatorname{gr}}
\newcommand{\pf}{\operatorname{pf}}
\newcommand{\Ext}{\operatorname{Ext}}
%\newcommand{\sup}{\operatorname{Tor}}
\newcommand{\Tor}{\operatorname{sup}}
\newcommand{\InjDim}{\operatorname{inj.dim}}
\newcommand{\MaxSpec}{\operatorname{MaxSpec}}
\newcommand{\Supp}{\operatorname{Supp}}
\newcommand{\Ker}{\operatorname{Ker}}
\newcommand{\Length}{\operatorname{Length}}
\newcommand{\CoKer}{\operatorname{Coker}}
\newcommand{\Max}{\operatorname{Max}}
\newcommand{\Ann}{\operatorname{Ann}}
\newcommand{\Ass}{\operatorname{Ass}}	
%\newcommand{\max}{\operatorname{max}}	
\DeclareMathOperator{\chara}{{{char}}}

\newcommand{\rank}{\operatorname{rank}}	
\newcommand{\bh}{\operatorname{bigheight}}	
\newcommand{\IM}{\operatorname{Im}}	
\newcommand{\depth}{\operatorname{depth}}	
\newcommand{\Char}{\operatorname{char}}
\newcommand{\GR}{\operatorname{gr}}
\newcommand{\Min}{\operatorname{min}}
\newcommand{\Proj}{\operatorname{Proj}}
\newcommand{\height}{\operatorname{height}}	
\newcommand{\rk}{\operatorname{rk}}	
%\newcommand{\rank}{\operatorname{rank}}	
\newcommand{\Ht}{\operatorname{ht}}	
\newcommand{\FITM}{\Pi-\hbox{modules}}
\newcommand{\FIN}{\Pi-\hbox{jumping number}}
\newcommand{\ceil}[1]{\lceil {#1} \rceil}
\newcommand{\deta}[1]{\operatorname{det}\left( {#1} \right)}
\newcommand{\floor}[1]{\lfloor {#1} \rfloor}
\newcommand{\FDer}[1]{\stackrel{#1}{\to}}
\newcommand{\FDerL}[1]{\stackrel{#1}{\longrightarrow }}
\newcommand{\onto}{\twoheadrightarrow}

\makeatletter
\@namedef{subjclassname@2020}{\textup{2020} Mathematics Subject Classification}
\makeatother


\newcommand{\up}[1]{\left\lceil #1 \right\rceil}
\newcommand{\down}[1]{\left\lfloor #1 \right\rfloor}
\newcommand{\tr}[2]{\left \langle {#1} \right \rangle_{#2}} 
\newcommand{\tail}[2]{ \left \ldbrack {#1} \right \rdbrack_{#2}}
\newcommand{\digit}[2]{ {#1}^{(#2)}}
\newcommand{\base}{\operatorname{base}}

\def \cJ{\mathcal J}
\def \cI{\mathcal I}
\def \cH{\mathcal H}
\def \R{\mathcal R}
\def \G{\mathcal G}
\def \N{\mathcal N}
\newcommand{\HH}[3]{\operatorname{H}^{#1}_{#2}\left(#3\right)}
%\newcommand{\ul}[1]{\underline{#1}}
\newcommand{\ol}[1]{\overline{#1}}
\newcommand{\tHH}[2]{\tilde{\operatorname{H}}_{#1}(#2)}


\newcommand{\tens}{\otimes}
\newcommand{\hhtest}[1]{\tau ( #1 )}
\renewcommand{\hom}[3]{\operatorname{Hom}_{#1} ( #2, #3 )}


\newcommand{\newLyu}{\widetilde{\lambda}}
\newcommand{\NewLyu}{\widetilde{\Lambda}}

%Connecting map in LES
\newcommand{\connarrow}{\ar `[d] `[l] `[llld] `[ld] [lld]}

%For notes
\newcommand{\note}[2]{ {\color{blue}
\vspace{2 mm}
\par \noindent
\reversemarginpar{\color{blue} \textbf{#1}}
\framebox{\begin{minipage}[c]{0.95 \textwidth}
 #2\end{minipage}}\vspace{5 mm}\par}}     

 %For red comments on typos
\renewcommand{\!}[1]{{\color{red}\text{\Huge$\star$ }#1}}

\DeclareMathOperator{\grade}{{grade }}
\DeclareMathOperator{\lcm}{{lcm}}

%Alessandro's commands
\newcommand{\ls}{\leqslant}%
\newcommand{\gs}{\geqslant}
\newcommand{\ds}{\displaystyle}
\newcommand{\eee}[1]{\underline{}^e \hspace{-0.05cm}{#1}}
\newcommand{\eeep}[1]{\underline{}^{e+e'} \hspace{-0.05cm}{#1}}
\newcommand{\p}{\mathfrak{p}}
%\newcommand{\q}{\mathfrak{q}}
\newcommand{\MIN}{\operatorname{Min}}
\newcommand{\ov}[1]{\overline{#1}}
\newcommand{\ann}{\operatorname{ann}}
\newcommand{\ck}[1]{{#1}^{\vee}}
\newcommand{\Q}{\mathbb{Q}}
\newcommand{\ehk}{\operatorname{e_{HK}}}
\newcommand{\ps}[1]{\llbracket {#1} \rrbracket}
\newcommand{\cd}{\operatorname{cd}}
\newcommand{\lct}{\operatorname{lct}}
\newcommand{\grc}{c_{\gr}}
\renewcommand{\a}{\mathfrak{a}}
\newcommand{\OO}{\mathcal{O}}
\newcommand{\JJ}{\mathcal{J}}
\newcommand{\M}{\mathcal{M}}
\newcommand{\C}{\mathcal{C}}
\newcommand{\II}{\mathbb{I}}
\newcommand{\init}{\mathrm{in}}
\newcommand{\ttt}{{\bf t}}
\DeclareMathOperator{\Sym}{Sym}
%\pagestyle{plain}
%-------------------------------------------------------


\makeatletter
\@addtoreset{equation}{section}
\makeatother
\renewcommand{\theequation}{\thesection.\arabic{equation}}



\author[A. De Stefani]{Alessandro De Stefani$^1$}
\address{Dipartimento di Matematica, Universit{\`a} degli Studi di Genova, Via Dodecaneso 35, 16146 Genova, Italy}
\email{destefani@dima.unige.it}
\thanks{$^{1}$ The first, fourth and fifth authors were partially supported by the PRIN 2020 project 2020355B8Y ``Squarefree Gröbner degenerations, special varieties and related topics''.}

\author[J. Montaño]{Jonathan Monta\~no$^2$}
\address{School of Mathematical and Statistical Sciences, Arizona State University, P.O. Box 871804, Tempe, AZ 85287-18041}
\email{montano@asu.edu}
\thanks{$^{2}$ The second author was partially funded by NSF Grant DMS \#2001645/2303605.}




\author[L. N{\'u}{\~n}ez-Betancourt]{Luis N{\'u}{\~n}ez-Betancourt$^3$}
\address{Centro de Investigaci{\'o}n en Matem{\'a}ticas, Guanajuato, Gto., M{\'e}xico}
\email{luisnub@cimat.mx}
\thanks{$^{3}$ The third author is  supported by CONACyT Grant \#284598.}
%\thanks{$^{**}$ The third author was partially supported by  CONACyT Grant 284598 and C\'atedras Marcos Moshinsky.}


\author[L. Seccia]{Lisa Seccia$^1$}
\address{Max Planck Institute for Mathematics in the Sciences, Leipzig, Germany}
\email{seccia@mis.mpg.de}

\author[M. Varbaro]{Matteo Varbaro$^1$}
\address{Dipartimento di Matematica, Universit{\`a} degli Studi di Genova, Via Dodecaneso 35, 16146 Genova, Italy}
\email{varbaro@dima.unige.it}

\subjclass[2020]{}

\keywords{}


\begin{document}



%\title[On the $F$-purity of affine Schubert and Ladder Determinantal Varieties ]{$F$-singularities of affine  Schubert and Ladder Determinantal Varieties }

\title[Ladder determinantal varieties and their symbolic blowups]% and affine schubert varieties ]
{Ladder determinantal varieties and their symbolic blowups} %and affine schubert varieties }



\begin{abstract} 
	 In this article we show that  the symbolic Rees algebra of a mixed ladder determinantal ideal is strongly $F$-regular.
	 Furthermore, we prove that the  symbolic associated graded algebra of a  mixed ladder determinantal ideal is $F$-pure. The latter implies that  mixed ladder determinantal rings are $F$-pure. We also show that  ideals  of the poset of minors of a generic matrix give rise to $F$-pure algebras with straightening law. 
	 %We also show that ideals defining alternating sign matrix varieties and ideals of poset of minors of a generic matrix give rise to $F$-pure algebras. %  with straightening law and recover  
\end{abstract}


\keywords{Ladder determinantal ideals,  symbolic powers, blowup algebras,  $F$-purity, $F$-regularity, Gr\"obner degenerations, matrix Schubert varieties, Knutson ideals.}
\subjclass[2020]{Primary: 14M12, 13C40, 13A35, 13A30. Secondary:  14M15.}


\maketitle
%\setcounter{secnumdepth}{1}
%\setcounter{tocdepth}{1}
%\tableofcontents

%%%%%%%%%%%%%%%%%%%%%%%%%%%%%%%%%%%%%%%%%%%%%%%%%%%%%%%%%
\section{Introduction}\label{Intro}
%%%%%%%%%%%%%%%%%%%%%%%%%%%%%%%%%%%%%%%%%%%%%%%%%%%%%%%%%
Let $X=(x_{i,j})$ be a $k\times \ell$ generic matrix. 
%Thinking at $X$ as the set $\{x_{ij}:(i,j)\in\{1,\ldots ,k\}\times \{1,\ldots ,\ell\}\}$,  
A subset $L\subseteq X$ is a  (two sided) ladder if whenever  $x_{i,j},x_{i',j'} \in L$ with $i \leq i'$ and $j \geq j'$,  we have $x_{u,v} \in L$ for every $i \leq  u \leq  i'$ and $j' \leq  v \leq  j$.   A ladder $L$ is  one-sided if the three `corners' $x_{1,1}, x_{1,\ell}, x_{k,\ell}$ belong to $L$. Let  $\kk$ be a field. The {\it {\rm (}unmixed\,{\rm )} ladder determinantal ideal} $I_t(L)$ is the ideal of $\kk[L]$ generated by the $t$-minors of $L$. The more general class of {\it mixed ladder determinantal ideals}, on the other hand, allows the sizes of the minors to vary  under  certain constraints,  %  according to a suitable recipe. 
%If the sizes are given by $\ttt=(t_1,\ldots , t_v)\in \ZZ_{>0}^v$, we denote the ideal by  $I_{\ttt}(L)\subseteq \kk[L]$    and when $t_1=t_2=\ldots =t_v=t$ we simply write  $I_{\ttt}(L)=I_t(L)$; see 
%These ideals, whose precise definition we give in
 see \autoref{sub_ladder} for precise definitions. The varieties determined by ladder determinantal ideals are called {\it ladder determinantal varieties.}

%were introduced by Gorla 
%\cite{GorlaMixed}. 

Ladder determinantal rings $\kk[L]/I_t(L)$ were introduced %in 1988 
by Abhyankar  in the study of singularities of Schubert varieties % of flag manifolds 
\cite{abhyankar1988enumerative}. 
Conca and Herzog \cite{concaherzog1997} proved that these rings have rational singularities if $\kk$ has characteristic zero; equivalently, ladder determinantal rings are $F$-rational if $\kk$ has characteristic $p\gg 0$ by the results of Smith and Hara \cite{SmithFrational, H98}. It is known that the special class of  chain ladders, which includes the one-sided ladders, are strongly $F$-regular in every positive characteristic \cite{concaherzog1997}, see also \cite{mehta1985frobenius}. Previously, Glassbrener and Smith had proved that complete intersection ladder determinantal rings are strongly $F$-regular \cite{GS}.  Nevertheless, the question of whether every  ladder determinantal ring  is strongly $F$-regular  remains open. In fact, up to now, also whether every ladder determinantal ring is $F$-pure has been an open question (see \cite[p. 122]{concaherzog1997}). In this article we prove the latter statement  more generally for mixed ladder determinantal rings.  We also show  the following  results over any field of characteristic $p>0$: % together the following results :
\begin{itemize}
	\item The symbolic associated graded ring of a mixed ladder determinantal ideal is $F$-pure (\autoref{ThmMixLadderSFP}, \autoref{blowups_f_pure}). In particular, mixed ladder determinantal rings are $F$-pure.
	\item The symbolic Rees algebra of a mixed ladder determinantal ideal is strongly $F$-regular (\autoref{stronglyF}).
\end{itemize}
%The reader should be aware that the notions of  $F$-purity and strongly $F$-regularity naturally extend to non-Noetherian rings; indeed we do not know whether  the symbolic blowup algebras above are Noetherian for every mixed ladder determinantal ideal.  
The Noetherianity of these algebras is known in the case unmixed ladders  by  work of Bruns and Conca  \cite[Theorem 4.1]{brunsconca98}, who proved the stronger statement that they admit a finite SAGBI basis with respect to an  antidiagonal term order. Using a combinatorial operation on ladders, which we call  {\it chamfering}, we show that they are Noetherian also in the case of mixed ladders (see \autoref{ThmNoetherian}). While our results do not guarantee that such algebras still admit a finite SAGBI basis, when they do we obtain additional results:
\begin{itemize}
	\item If the symbolic Rees algebra of a mixed ladder determinantal ideal $I$ admits a finite SAGBI basis with respect to an antidiagonal term order and $\chara(\kk)\gg 0$, then $\IN\left(I^{(n)}\right)=\IN(I)^{(n)}$ for all $n\in\NN$ (\autoref{cor_SAGBI}).
\end{itemize}

We give a second proof that unmixed ladder determinantal rings are $F$-pure, proving that their defining ideals are {\it Knutson}  associated to the product $f$ of the minors corresponding to  the main antidiagonals of the ladder. The notion of Knutson ideals  was  introduced   over finite fields  \cite{knutson2009frobenius} and later developed over any field (also of characteristic 0) by Seccia \cite{seccia2021knutson}. More generally we prove:
\begin{itemize}
	\item Given a ladder, the ideal generated by fixed-size minors in the subladders corresponding to consecutive rows or columns in the matrix are Knutson ideals associated to the  polynomial $f$ described above. In particular, the natural minimal generators of any sum of these ideals form a Gr\"obner basis with respect to an antidiagonal term order. If, besides the operation of taking sums, we also allow the one of taking intersections, all the resulting ideals still admit a squarefree initial ideal with respect to an antidiagonal term order and, in positive characteristic, define $F$-pure rings (\autoref{ThmLadderKI}). 
\end{itemize}
When $L=X$ the above result had already been proved by Seccia  \cite{seccia2022knutson}. In the case $L=X$ we also prove that any ideal of fixed-size minors of a top-left or bottom-right corner submatrix of $X$ is Knutson associated to $f$ (\autoref{p:msk}). As a consequence any sum of these ideals is  Knutson associated to $f$,  their natural minimal generators  form a Gr\"obner basis with respect  to antidiagonal term orders. If, besides the operation of taking sums, we also allow the one of taking intersections, all the resulting ideals still admit a squarefree initial ideal with respect to an antidiagonal term order and, in positive characteristic, define $F$-pure rings. Some classes of ideals that can be obtained this way and for which we obtain the latter consequences  are: %From this fact we get two consequences:
\begin{itemize}
	\item Schubert determinantal ideals, and more generally, alternating sign matrix ideals (as they are sums of ideals of minors of top-left corner submatrices of $X$), see %is a Knutson ideal for $f$ and its natural minimal generators form a Gr\"obner basis w.r.t. an antidiagonal term order. In particular, in positive characteristic matrix Schubert varieties are $F$-split 
	\autoref{c:msk}.
	\item Any ideal of the poset of minors of $X$ (as any of them is the intersection of sums of ideals of minors of top-left corner submatrices of $X$), see  % generates a Knutson ideal for $f$. In particular, in positive characteristic any ASL coming from an ideal of the poset of minors is $F$-pure 
	\autoref{c:pif}. %Thus, in positive characteristic any algebra with straightening law (ASL) arising from an ideal of the poset of minors is $F$-pure. 
\end{itemize}

We remark that Schubert determinantal ideals % matrix Schubert varieties 
were already known to be  %The results described in the first of the previous points were already known: that matrix Schubert varieties are 
%globally 
$F$-pure in positive characteristic by work of Brion and Kumar  \cite{Brion_Kumar_Book}. The fact that the natural minimal generators of Schubert determinantal ideals form a Gr\"obner basis with respect  to antidiagonal term orders was also known  by work of Knutson and Miller  \cite{knutson2005grobner}. Moreover,  a proof by induction on the Bruhat order, and so different from the one given here, of the fact that  Schubert determinantal ideals are Knutson associated to $f$ was  already sketched  in \cite{knutson2009frobenius}. Klein and  Weigandt \cite{klein2021bumpless} showed that alternating sign matrix ideals have radical initial ideals with respect to antidiagonal term orders and hence they are radical themselves. 
All of the above facts are recovered, and sometimes strengthened, by \autoref{c:msk}. 

On the other hand, regarding \autoref{c:pif}, we highlight the connection with  Eisenbud's suspicion that graded ASLs might always be $F$-pure in positive characteristic \cite[p. 245]{eisenbud80}. As noticed by Koley and Varbaro \cite[Remark 5.2]{KoleyVarbaro} this is not always true, although  all graded ASLs are $F$-injective. \autoref{c:pif} confirms Eisenbud's intuition for ASLs arising from ideals of the poset of minors of a generic matrix. 

%% talk about symbolic powers 

%%%%%%%%%%%%%%%%%%%%%%%%%%%%%%%%%%%%%%%%%%%%%%%%%%%%%%%%%
\section{Background}\label{SectionBackground}
%%%%%%%%%%%%%%%%%%%%%%%%%%%%%%%%%%%%%%%%%%%%%%%%%%%%%%%%%

In this section we include some preliminary material that is needed in the subsequent sections. In particular, in   \autoref{sub_ladder} and \autoref{sub_affine_schub} we recall the definitions of the classes of determinatal ideals that are treated in this paper. 

%%%%%%%%%%%%%%%%%%%%%%%%%%%%%%%%%%%%%%%%%%%%%%%%%%%%%%%%%
\subsection{Methods in prime characteristic}\label{SubSecPrimeChar} %%%%%%%%%%%%%%%%%%%%%%%%%%%%%%%%%%%%%%%%%%%%%%%%%%%%%%%%%

We begin by recalling some concepts from positive characterisic commutative algebra. For more infomation we refer the reader to Ma and Polstra's recent survey \cite{ma2021f}.

\begin{definition}\label{DefFbasic}
Let $R$ be a ring of prime characteristic $p$. 
%\begin{itemize}[]
%\item 

For every $e\in \ZZ_{> 0}$, we denote by $F^e_*( R)=\{F^e_* (r)\;|\; r\in R\}$ the $R$-algebra which is isomorphic to $R$ as a ring and whose action by $R$ is given by $f F^e_* (r)=F^e_* \left(f^{p^e}r\right)=F^e_*\left(f^{p^e}\right) F^e_* (r)$. % If $e=1$. we simply write $F^e_*$  as $F_*$.  
We say that $R$ is {\it $F$-finite} if the the Frobenius map $F^e:R\to F_*^e(R)$, $r\mapsto F_*^e(r^{p^e})$ is finite for some, or equivalently all, $e\in \ZZ_{> 0}$, that is, if  $F_*^e(R)$ is a finitely generated  $R$-module. 
%\item 
%
%\item 

We say that $R$ is  {\it $F$-pure} if for some, or equivalently all, $e\in \ZZ_{>0}$ the Frobenius map $F^e$ is pure, i.e.,  the map $M\otimes_R R\xrightarrow{\mathbf{1}_M\otimes_R F^e} M\otimes_R F_*^e(R)$ is injective for every $R$-module $M$. 
%\item 
We say that $R$ is {\it $F$-split} if for some, or equivalently all, $e\in \ZZ_{>0}$ the map $ F^e$ splits. 
Clearly, if $R$ is $F$-split, it is $F$-pure; the converse holds if $R$ is $F$-finite. We also note that if $R$ is  $F$-pure, then it is reduced. Therefore, $R$ is $F$-pure (resp. $F$-split) if and only if the map $R\to R^{1/p^e}$ is pure (resp. splits) for   some, or equivalently all, $e\in \ZZ_{>0}$. 
%\end{itemize}
\end{definition}





We now  recall the following notion that connects symbolic powers and the property of being $F$-split \cite{DSMNB}. 

\begin{definition}\label{DefFpureFilt}
 Let $R=\kk [x_1,\ldots,x_d]$ be a standard graded polynomial ring over an $F$-finite  field $\kk$ of prime characteristic. We say that a homogeneous ideal $I\subseteq R$ is  {\it symbolic $F$-split} if  there exists a splitting $\phi:R^{1/p}\to R$ such that
$\phi\left(\left(I^{(np+1)}\right)^{1/p}\right)\subseteq I^{(n+1)}$ for every $n\in \NN$.   
Equivalently, if there exists 
$$f\in \bigcap^{\infty}_{n=1} \left( \left( I^{(n+1)}\right)^{[p]}:I^{(np+1)}\right)\setminus \m^{[p]}.$$
\end{definition}

\begin{remark}\label{rem_impt_prop_symb}
	In the present paper, ideals which are symbolic $F$-split are especially relevant because of the following two properties:
\begin{enumerate}
\item If  $I\subseteq R$ is  symbolic $F$-split, 
then the {\it symbolic Rees algebra} $\R^s(I):=\oplus_{n\in \NN} I^{(n)}$ and the {\it symbolic associated graded algebra}   $\gr^s(I):=\oplus_{n\in \NN} I^{(n)}/I^{(n+1)}$ are $F$-split. 
In particular, if $I$ is symbolic $F$-split then  $R/I$ is also $F$-split \cite[Theorem 4.7]{DSMNB}.
\item Let  $I\subseteq R$ be a homogeneous ideal and $H:=\bh(I)$ the maximum height of a minimal prime ideal of $I$. IF there exists a term order $\prec$ on $R$ for which $\init_\prec\left(I^{(H)}\right)$ contains a squarefree monomial, then $I$ is symbolic $F$-split  \cite[Lemma 6.2(1)]{DSMNB}.  Moreover, in this case $\init_\prec(I)$ is automatically radical \cite[Theorem 3.13]{KoleyVarbaro},  and if $I$ is equidimensional then both $\bigoplus_{n\in \NN}\init_\prec(I^{(n)})$ and $\bigoplus_{n\in \NN}\frac{\init_\prec(I^{(n)})}{\init_\prec\left(I^{(n+1)}\right)}$ are $F$-split \cite[Proposition 7.5]{DSMNB}.% We note that the assumption in  \autoref{DefFpureFilt}(2) automatically implies that . %; in fact, this holds  even  in characteristic zero. 
\end{enumerate}
\end{remark}

%There is a stronger relationship between symbolic powers and initial ideals in this context. To be able to discuss it,  let us first recall that,
Let $I\subseteq R$ be a homogeneous ideal such that  $\init_\prec(I)$ is radical, then   
$\init_\prec\left (I^{(n)}\right)\subseteq \init_\prec(I)^{(n)}$  for every  $n\in\NN$   \cite[Proposition 5.1]{Sull}. 
%On the other hand, even if $\init_\prec(I)$ is not radical, we always have 
%$$\init_\prec(I)^n\subseteq \init_\prec\left (I^n\right)\subseteq \init_\prec\left (I^{(n)}\right) \text{ for every } n\in\NN.$$ Hence, when $\init_\prec(I)$ is radical, we have the inclusions
Hence, \[\init_\prec(I)^n\subseteq \init_\prec\left (I^{(n)}\right)\subseteq \init_\prec(I)^{(n)} \,\,\,\text{ for every }n\in\NN.\]
Therefore, since $\init_\prec(I)$ is radical,  $\init_\prec\left (I^{(n)}\right)=\init_\prec(I)^{(n)}$ if and only if $\init_\prec\left (I^{(n)}\right)$ has no embedded components. This fact allows us to show the following interesting relationship between symbolic powers and initial ideals.
%
%\begin{proposition}\label{symbolic-initial}
%Assume $\kk$ is a perfect field and let $I\subseteq R=\kk[x_1,\ldots ,x_d]$ be a homogeneous equidimensional radical ideal of height $h$. Fix a monomial order $\prec$ on $R$ and  assume that $\bigoplus_{n\in \NN}\init_\prec(I^{(n)})$ is Noetherian. If $\chara(\kk)\gg 0$, the following are equivalent
%\begin{enumerate}[\rm (1)]
%\item $\init_\prec\left(I^{(n)}\right)=\init_\prec(I)^{(n)}$ for every $n\in\NN$ and $\init_\prec(I)$ is radical.
%\item $\init_\prec\left(I^{(h)}\right)=\init_\prec(I)^{(h)}$ and $\init_\prec(I)$ is radical.
%\item $\init_\prec\left(I^{(h)}\right)$ contains a squarefree monomial.
%\end{enumerate}
%\end{proposition}
%\begin{proof}
%The implication $(1)\Rightarrow (2)$ is trivial. For $(2)\Rightarrow (3)$ we note that  $x_1\cdots x_d$ belongs to $P^h$ for every minimal prime ideal $P$ of $\init_\prec(I)$ and then $x_1\cdots x_d\in \init_\prec(I)^{(h)}= \init_\prec\left(I^{(h)}\right)$. 
%It remains to show  $(3)\Rightarrow (1)$.  From \autoref{rem_impt_prop_symb}(2) it follows that  $\IN_\prec(I)$ is radical and $\bigoplus_{n\in\NN}\frac{\init_\prec(I^{(n)})}{\init_\prec(I^{(n+1)})}$ is reduced. Thus, $\init_\prec\left(I^{(n)}\right)=\init_\prec(I)^{(n)}$ for every $n\in\NN$  \cite[Corollary 7.10]{DSMNB} \jonathan{Don't we need $I$ prime here?}.
%\end{proof}


\begin{proposition}\label{symbolic-initial}
	 Let $R=\kk [x_1,\ldots,x_d]$ be a standard graded polynomial ring over a perfect field $\kk$ of prime characteristic. 
%	Assume $\kk$ is a perfect field and let 
Let $P\subseteq R$ be a homogeneous prime ideal of height $h$. Fix a term order $\prec$ on $R$ and  assume that $\bigoplus_{n\in \NN}\init_\prec\left(P^{(n)}\right)$ is Noetherian. If $\chara(\kk)\gg 0$, the following are equivalent
	\begin{enumerate}[\rm (1)]
		\item $\init_\prec\left(P^{(n)}\right)=\init_\prec(P)^{(n)}$ for every $n\in\NN$ and $\init_\prec(P)$ is radical.
		\item $\init_\prec\left(P^{(h)}\right)=\init_\prec(P)^{(h)}$ and $\init_\prec(P)$ is radical.
		\item $\init_\prec\left(P^{(h)}\right)$ contains a squarefree monomial.
	\end{enumerate}
\end{proposition}
\begin{proof}
	The implication $(1)\Rightarrow (2)$ is trivial. For $(2)\Rightarrow (3)$ we note that  $x_1\cdots x_d$ belongs to $Q^h$  for every minimal prime ideal $Q$ of $\init_\prec(P)$ and then $x_1\cdots x_d\in \init_\prec(P)^{(h)}= \init_\prec\left(P^{(h)}\right)$. 
	It remains to show  $(3)\Rightarrow (1)$.  From \autoref{rem_impt_prop_symb}(2) it follows that  $\IN_\prec(P)$ is radical and $\bigoplus_{n\in\NN}\frac{\init_\prec(P^{(n)})}{\init_\prec(P^{(n+1)})}$ is reduced. Thus, $\init_\prec\left(P^{(n)}\right)=\init_\prec(P)^{(n)}$ for every $n\in\NN$  \cite[Corollary 7.10]{DSMNB}. 
\end{proof}

\begin{remark}
The lower bound on the characteristic of the field $\kk $ in \autoref{symbolic-initial} can be made precise: if $\bigoplus_{n\in \NN}\init_\prec(P^{(n)})$ is generated as an $R$-algebra by forms of  degrees $b_1,\ldots ,b_m$, then \autoref{symbolic-initial} applies if  $\chara(\kk)>\lcm(b_1,\ldots ,b_m)$.
\end{remark}





%%%%%%%%%%%%%%%%%%%%%%%%%%%%%%%%%%%%%%%%%%%%%%%%%%%%%%%%%
\subsection{Knutson ideals}\label{SubSecKnutson}
%%%%%%%%%%%%%%%%%%%%%%%%%%%%%%%%%%%%%%%%%%%%%%%%%%%%%%%%%

The notion of {\it Knutson ideals} was coined by  Conca and Varbaro \cite{conca2020square}  after Knutson's work on compatibly split ideals and degeneration 
\cite{knutson2009frobenius} (see also \cite{brion2007frobenius}).

\begin{definition}\label{K.I.} Let $f \in R= \kk [x_1,\ldots,x_d]$ be a homogeneous polynomial where $\kk$ is a field and assume that  for some term order $\prec$ the leading term of $f$ is a squarefree monomial. 
	We define $\mathcal{C}_f$ to be the smallest set of ideals satisfying the following conditions:
	\begin{enumerate}[(1)]
		\item $(f) \in \mathcal{C}_f$.
		\item  If $I \in \mathcal{C}_f$ then $(I:J) \in \mathcal{C}_f$ for every ideal $J \subseteq R$.
		\item If $I$ and $J$ are in $\mathcal{C}_f$ then also $I+J$ and $I \cap J$ must be in $\mathcal{C}_f$.
	\end{enumerate} 
	If $I$ is an ideal in $\mathcal{C}_f$, we say that $I$ is a \emph{Knutson ideal associated to} $f$, or simply that it is {\it Knutson}. % . More generally, we say that $I$ is a \emph{Knutson ideal} if $I \in \mathcal{C}_f $ for some $f$.
\end{definition} 
If $\kk$ has positive chacteristic and $\deg(f)=d$, then $f$ defines a splitting map on $R$  (see \cite[Lemma 4]{knutson2009frobenius} for finite fields and  \cite{seccia2021knutson} in general). %Therefore,
Knutson ideals are %turn out to be 
compatibly split with respect to this map. Theorefore,  Knutson ideals  define $F$-split %$F$-pure (equivalently, $F$-split)
rings in positive characteristic. In addition, Knutson ideals
have squarefree initial ideals  in arbitrary characteristic (see \cite[Theorem 2]{knutson2009frobenius} for finite fields and \cite{seccia2021knutson} in general) and consequently their extremal Betti numbers %, Castelnuovo-Mumford regularity, and depth of Knutson ideals 
coincide
with  those of their initial ideals \cite{conca2020square}. Furthermore, Gr\"obner bases of Knutson ideals are well-behaved %\lq\lq well\rq \rq  
with respect to sums as the %in the sense  that the 
union of Gr\"obner bases of Knutson ideals is a Gr\"obner basis of their sum. This quality makes computations of Gr\"obner bases for some Knutson ideals somewhat  easier than for general ideals. % and then it makes  them useful for solving problems in applied algebra. %(see for example \cite{}\jonathan{Add reference}). %\par 
All of  these considerations make  Knutson ideals relevant  objects %f particular interest 
in computational algebra, combinatorial commutative algebra, and $F$-singularity theory.  
%We refer the reader to Seccia's  recent papers  for examples and further  details about the properties of Knutson ideals in any characteristic \cite{seccia2021knutson, seccia2022knutson}. 



\begin{remark}
	Since every ideal of $\mathcal{C}_f$ is radical,  condition   \autoref{K.I.}(2) can be replaced by:
	\begin{itemize}
		\item[$(2^\prime)$] If $I \in \mathcal{C}_f$ then $P \in \mathcal{C}_f$ for every $P \in \MIN(I)$.
	\end{itemize}
\end{remark}
%More details about the properties of Knutson ideals in any characteristic can be found in \cite{seccia2021knutson}. These properties were first proved by Knutson \cite{knutson2009frobenius} in the case $\kk=\mathbb{Z}/p\mathbb{Z}$. 



%%%%%%%%%%%%%%%%%%%%%%%%%%%%%%%%%%%%%%%%%%%%%%%%%%%%%%%%%
\subsection{Ladder determinantal ideals}\label{sub_ladder}
%%%%%%%%%%%%%%%%%%%%%%%%%%%%%%%%%%%%%%%%%%%%%%%%%%%%%%%%%


In this subsection we recall  the definition of mixed ladder determinantal ideals. For this, we follow closely the notations and conventions form Gorla's work \cite{GorlaMixed}, where these ideals were introduced    (see also \cite{gonciulea2000mixed}).

\begin{definition}\label{def_ladder}
	Let $X=(x_{i,j})$ be a generic matrix of size $k \times \ell$,  i.e.,
	\[ X=
	\begin{bmatrix}
		x_{1,1}       & \dots  & x_{1,\ell} \\
		%	x_{2,1}        & \dots  & x_{2,\ell} \\
		\vdots  & \vdots &\vdots \\
		%   x_{k-1,1}       & x_{k-1,2}  & \dots & x_{k-1,l}\\
		x_{k,1}        & \dots & x_{k,\ell}
	\end{bmatrix}.
	\]% We index the entries of $X$ following the linear algebra convention, i.e., the entry $(i,j)$ is the one in row $i$ and column $j$.  
	%	Let $X=(x_{ij})$ be a $k \times \ell$ matrix of indeterminates and $L\subseteq  X$. 
	We say that a subset $L\subseteq X$ is a {\it {\rm (}two sided\,{\rm )} ladder} if it has
	the following property: if $x_{i,j},x_{i',j'} \in L$ with $i \leq i'$ and $j \geq j'$,  then $x_{u,v} \in L$ for any $i \leq  u \leq  i'$ and $j' \leq  v \leq  j$.  
	%
	%Let $L$ be a ladder. 
	We denote by 
	$$\left\{x_{1,a_1}=x_{b_1,a_1},\ldots,x_{b_u,a_u}=x_{b_u,\ell}\right\} \text{ with } 1=b_1<\cdots < b_u\text{  and  } a_1<\cdots < a_u=\ell$$ 
	%$$\left\{(1,a_1)=(b_1,a_1),\ldots,(b_u,a_u)=(b_u, \ell)\right\} \text{ with } 1=b_1<\cdots < b_u$$ 
	the {\it upper-outside corners}, and by 
	$$\left\{x_{d_1,1}=x_{d_1,c_1},\ldots,x_{d_v,c_v}=x_{k,c_v}\right\} \text{ with } d_1\ls \cdots \ls d_v=k
	\text{  and } 1=c_1\ls \cdots \ls c_v$$ the {\it lower-outside corners} of $L$.  Moreover, we assume  all of the corners are distinct. %and that and no two consecutive corners coincide and that 
	%$b_i < b_j$ and $c_i < c_j$ if $i < j$.
	If $L$ has only one upper-outside corner we say it is {\it one-sided}. 
	We also define the following subladders for all $j=1,\ldots ,v$
\begin{equation}\label{eq_Lj}
	L_j = \{x_{b,a} \in L \mid 1 \leq b \leq d_j, c_j \leq a \leq l\}.
\end{equation}
%	Note that each $L_j$  has a %is a one-sided ladder, i.e. with only one lower- or upper-outside corner. We note that the  
	%with 
%	unique lower-outside corner  at %of $L_j$ is 
%	$x_{d_j,c_j}$, such a ladder is said to be  {\it one-sided}. %and $L=\bigcup_{j=1}^v L_j$. 
	In \autoref{example_ladder_1} we illustrate all of these notations in a specific example. 
\end{definition}


\begin{example}\label{example_ladder_1}
	Let $X$ be a generic matrix of size $10\times 10$. In \autoref{ladder} the shaded area represents the ladder $L$ whose upper-outside corners are $\{x_{1,2},\,x_{4,8},\, x_{8,10}\}$ and lower-outside corners are $\{x_{3,1},\,x_{6,1},\, x_{8,6},\, x_{10,8}\}$; note that the first two lower-outside corner belong to the same column. 
	\begin{figure}[ht]
		\includegraphics[scale=0.4]{Ladder_whole.eps}
		\caption{Ladder  $L$}
		\label{ladder}
	\end{figure}
	
	\noindent Since $L$ has four lower-outside corners, it has four subladders $L_1,L_2, L_3, L_4$ as defined above. 
	In \autoref{subladders} we highlight these subladders. % in this order from left to right. 
	\begin{figure}[ht]
		\centering
		\mbox{\subfigure{\includegraphics[width=1.5in]{sub_ladder_1.eps}}\quad
			\subfigure{\includegraphics[width=1.5in]{sub_ladder_2.eps} } \quad
			\subfigure{\includegraphics[width=1.5in]{sub_ladder_3.eps} }\quad
			\subfigure{\includegraphics[width=1.5in]{sub_ladder_4.eps} }
		}
		\caption{The subladders $L_1$, $L_2$, $L_3$, $L_4$ ordered from left to right} \label{subladders}
	\end{figure}
\end{example}


\begin{definition}\label{def_mixed_lad_det}  Let $\kk$ be a field and $L \subseteq X$  a ladder as in \autoref{def_ladder}. If  $t$ is a positive integer, we denote by $I_t (L)$ the ideal of $\kk[L]$ generated by the  $t$-minors of $X$ that are completely contained in $L$. The ideal $I_t (L)$ is a {\it ladder determinantal ideal} of $L$. On the other hand, if $\ttt=(t_1,\ldots,t_v) \in \ZZ_{>0}^v$ we denote by $I_\ttt(L)$ the ideal $I_{\ttt}(L) = I_{t_1}(L_1) + \ldots + I_{t_v}(L_v)$. Such ideal is a {\it mixed ladder determinantal ideal} of $L$. Notice that,  if $t=t_1=\cdots=t_v$, then  $I_{\ttt}(L)=I_t (L)$. The ideals $I_{\ttt}(L)$ and $I_{t}(L)$ are prime (see \cite[]{GorlaMixed}).
	
	Throughout, for every mixed ladder determinantal ideal  we assume that  all the entries of $L$ are involved in some generator of  $I_{\ttt}(L) $ and that  $I_{t_i}(L_i) \not\subseteq I_{t_j}(L_j)$ for all $i \ne j$.
\end{definition}

\begin{example}
	Let $L$ be as in \autoref{example_ladder_1}. If $\ttt=(2,3,1,2)$, then $I_\ttt(L)=I_2(L_1)+
	I_3(L_2)+
	I_1(L_3)+
	I_2(L_4).$  If $\ttt=(2,2,2,2)$, then $I_\ttt(L)=I_2(L)$.
\end{example}



\begin{notation} \label{running assumptions ladder} We summarize here the assumptions on the ladders treated in the present papers, which are the same as in  Gorla's work \cite{GorlaMixed}.
	\begin{enumerate}[(1)]
		\item All the entries of $L$ are involved in some minor, otherwise we delete the superfluous entries and consider the remaining  ladder. 
		\item No two consecutive corners coincide, and no two upper-outside corners belong to the same row or to the same column.
		\item  $I_{t_i}(L_i) \not\subseteq I_{t_j}(L_j)$ for all $i \ne j$
		\item When we view $L$ as a ladder inside the $k \times \ell$ matrix $X$ we assume that there is no proper submatrix $X'$ of $X$ which contains $L$.
	\end{enumerate}
\end{notation}


%For $1 \leq i \leq j \leq v$, we let 
%\[
%L_{i,j} = (L_i \smallsetminus L_{i-1}) \cap (L_j \smallsetminus L_{j+1}),
%\]
%where we set $L_{-1} = L_{v+1} = \emptyset$. Moreover, set
%\[
%t_{i,j} = \begin{cases} \min\{t_u \mid L_{i,j} \subseteq L_u\} & \text{ if } L_{i,j} \ne \emptyset \\  1 & \text{ otherwise }
	%\end{cases}
	%\]



%%%%%%%%%%%%%%%%%%%%%%%%%%%%%%%%%%%%%%%%%%%%%%%%%%%%%%%%%
\subsection{Schubert determinantal ideals and  poset ideals}\label{sub_affine_schub}
%%%%%%%%%%%%%%%%%%%%%%%%%%%%%%%%%%%%%%%%%%%%%%%%%%%%%%%%%


 
In this  subsection we introduce further families of determinantal ideals treated in this paper. %corresponding to affine Schubert varieties.  
We begin by introducing some notations. 

\begin{notation}\label{notation_affine_subsec}
%	Let $X=(x_{i,j})$ be a generic matrix of size $k \times \ell$, %and entries $x_{i,}$, 
%	i.e.,  
%	\[ X=
%	\begin{bmatrix}
%		x_{1,1}       & \dots  & x_{1,\ell} \\
	%	x_{2,1}        & \dots  & x_{2,\ell} \\
%		\vdots  & \vdots &\vdots \\
		%   x_{k-1,1}       & x_{k-1,2}  & \dots & x_{k-1,l}\\
%		x_{k,1}        & \dots & x_{k,\ell}
%	\end{bmatrix}.
%	\]
	Let $X$ be as in \autoref{def_ladder}.	For any $1 \leq a <b \leq \ell$ and $1 \leq c<d \leq k$, we denote by $X_{[a,b]}^{[c,d]}$ the submatrix of $X$ with column indices $a,\ldots, b$ and row indices $c,\ldots, d$. 
	In the case $[c,d]=[1,k]$ (resp. $[a,b]=[1,\ell]$ ) we simply write $X_{[a,b]}$ (resp. $X^{[c,d]}$).  
	We use $X_{a \times b}$ to denote the northwest submatrix $X_{[1,b]}^{[1,a]}$ and $X^{a \times b}$ to denote the southeast submatrix $X_{[\ell-b+1 ,\ell]}^{[k-a+1,k]}$. 
	 Moreover, for a set of indices $1\ls i_1<\cdots<i_r\ls k$ and $1\ls j_1<\cdots<j_r\ls \ell$ we denote by $[i_1\cdots i_r\mid j_1\cdots j_r]$ the determinant ($r$-minor) of the submatrix of $X$ corresponding to rows $i_1,\ldots, i_r$ and columns  $j_1,\ldots, j_r$.
\end{notation}

Following \cite{MillerSturmfels}, by the {\it antidiagonal term} of a minor $[i_1\cdots i_r\mid j_1\cdots j_r]$ we mean the product of the variables in the antidiagonal of the corresponding submatrix, namely $x_{i_1,j_r}x_{i_2,j_{r-1}}\cdots x_{i_r,j_1}$. Let $\kk$ be an arbitrary  field.  A term order on the polynomial ring $\kk[X]$ is {\it antidiagonal} if the initial term of any minor is  its antidiagonal term. A typical example of an antidiagonal term order is the lexicographical one with variables ordered as $$x_{1,\ell}>\cdots>x_{1,1}>x_{2,\ell}>\cdots>x_{2,1}>\cdots >x_{k,\ell}>\cdots>x_{k,1}.$$

Although it is not our intention to review the theory of {\it algebras with straightening law} in this paper (see \cite[Chapter 3]{bruns2022determinants}),  we wish to introduce a few concepts for the sake of clarity. The set $\Pi$ of all the minors of the generic matrix $X$ can be partially ordered by the relation
\[[a_1\ldots a_r|b_1\ldots b_r]\leq [c_1\ldots c_s|d_1\ldots d_s] \iff r\geq s, \ \ a_i\leq c_i \mbox{ and } b_i\leq d_i \ \text{ for all } \ i=1,\ldots ,s.\] 
 The starting observation is that $\kk[X]$  is generated %by $\Pi$ 
as a % $\kk$-algebra. In fact, a basis of $\kk[X]$ as 
$\kk$-vector space by the set %is given 
by
\[\{\pi_1\cdots \pi_m\mid  \ \pi_i\in\Pi \text{ for }i=1,\ldots, m \text{ and } \ \pi_1\leq \cdots \leq \pi_m\}.\]
The elements of this $\kk$-basis are called {\it standard monomials}. Althoug it may happen that the product of two standard monomials is not a standard monomial, such product will be uniquely writable as a $\kk$-linear combination of standard monomials;   %, which is in some sense compatible with the poset structure on $\Pi$.
 this process  is known in the literature as {\it straightening law}. 
 
With the above notation, we are ready to define matrix Schubert varieties and its defining ideals (see \cite[Chapter 15]{MillerSturmfels} for more information).

\begin{definition}\label{def_MSV}
%\begin{enumerate}[(1)]
Let $\mathscr{M}_{k\times \ell}(\mathbb{\kk})$ be the vector space of $k\times \ell$ matrices with entries in $\kk$ and let $w=(w_{i,j})\in \mathscr{M}_{k\times \ell}(\mathbb{\kk})$ be a {\it partial permutation}, i.e., $w_{i,j}=0$ for all $i,j$ except for at most one entry equal to 1 in each row and  column.
\begin{enumerate}[(1)]
\item The {\it matrix Schubert variety} $\overline{X_w}$ associated to $w$ is given by 
$$\overline{X_w}:=\left\{Y\in \mathscr{M}_{k\times \ell}\mid \rank(Y_{r\times s})\ls \rank(w_{r\times s}), \text{  for all }1\ls r\ls k, 1\ls s\ls \ell
\right\}.
$$
\item The defining ideal of $\overline{X_w}$ is the {\it Schubert determinantal ideal} $I_w$ which is generated by the $\rank(w_{r\times s})+1$ minors of $X_{r\times s}$ for every $1\ls r\ls k, 1\ls s\ls \ell$.
\end{enumerate}	

%\end{enumerate}
\end{definition}

\begin{remark}
	\
	\begin{enumerate}[(1)]
		\item Given $1\ls t\ls \min\{k,\ell\}$, if $w\in \mathscr{M}_{k\times \ell}(\mathbb{\kk})$ is the partial permutation defined by  $w_{ii}=1$ for $1\ls i\ls t-1$, then $I_w$ is the classical determinantal ideal of $t$-minors of $X$. 
			\item 
		Mixed ladder determinantal ideals of  ladders with only one upper-outside corner coincide with the Schubert determinantal ideals of {\it vexillary permutations} \cite[Proposition 9.6]{fulton1992flags}.
		\item For any partial permutation  there are smaller sets of generators for $I_w$ (see \cite[Theorem 15.15]{MillerSturmfels}, \cite{gao2022minimal}).

	%	\item If $k=\ell$, $\kk=\mathbb{C}$, and $w$ is a (full rank) permutation matrix, the varieties $\overline{X_w}$ are the ones mentioned in \autoref{Intro}. % of this subsection.  
	   
    \end{enumerate}
\end{remark}
%
% \textcolor{blue}{
%MSV:
%\begin{enumerate}
%\item Discuss determinantal ideals of MSV,  partial permutation, Rothe Diagram, Example 1  - PITTURE 1, and Example 2- PICTURE 2.
%
%\end{enumerate} 
%}

%\begin{example}
%	content...
%\end{example}

We continue with the definition of the following class of ideals.
\begin{definition}
	An {\it  ideal} of $\Pi$ is a subset $\Omega\subseteq \Pi$ such that $$\text{ for all } \omega\in \Omega \text{ and } \pi\in\Pi, \text{  if } \pi\leq \omega \text{ then } \pi\in\Omega.$$  
For any $\Omega$ ideal of $\Pi$, we call the ideal generated by $\Omega$ in $\kk[X]$  a {\it poset ideal}.
\end{definition}

%
%\textcolor{blue}{
%\begin{enumerate}
%\item  EXAMPLE.
%\item 
%\end{enumerate}
%}

\begin{remark}
A particular type of  ideals of $\Pi$, fixed $\delta\in\Pi$,  are the sets $$\Omega_\delta:=\{\pi\in\Pi: \pi\not\gs \delta\}.$$
The poset ideals $\Omega_\delta\kk[X]$ are %precisely 
the defining ideals of %the  
varieties %mentioned in  \autoref{Intro} 
originating as certain affine charts of Schubert varieties  \cite[p. 535]{procesi2007lie} (see also \cite[Chapter 5, \S A]{BookDet}).
\end{remark}



%%%%%%%%%%%%%%%%%%%%%%%%%%%%%%%%%%%%%%%%%%%%%%%%%%%%%%%%%%%
\section{Chamfering ladders}\label{Decorn_Section}
%%%%%%%%%%%%%%%%%%%%%%%%%%%%%%%%%%%%%%%%%%%%%%%%%%%%%%%%%%%
	
In this section we define a combinatorial operation on ladders, which we call  {\it chamfering}, and show some important consequences of it. This definition is inspired by the work of previous authors \cite{conca1995ladder, concaherzog1997, GorlaMixed}.

\begin{setup} \label{setup_chamfer} 
	We adopt the assumptions and notations from \autoref{def_ladder}, \autoref{def_mixed_lad_det}, and \autoref{running assumptions ladder}. % follow the definitions and notations from \autoref{sub_ladder}.
	In particular, $X=(x_{i,j})$ is a generic matrix of size $k \times \ell$ and $L\subseteq	 X$ is a ladder. Let $\kk$ be a field. We denote by $R$ the polynomial ring $R=\kk[L]$. Given   $\ttt=(t_1,\ldots,t_v) \in \ZZ_{>0}^v$ we denote by $I_\ttt(L)$ the corresponding mixed ladder determinantal ideal.    
\end{setup}


\begin{definition}\label{def_chamfer}
Assume \autoref{setup_chamfer}. We say a ladder $L$ is {\it non-degenerate} with respect to $\ttt$ if $I_\ttt(L)$ satisfies the assumptions in \autoref{running assumptions ladder}.  % Let $x_{d_j, c_j}$ be such that $d_{j-1}<d_j$ and $c_j>c_{j+1}$. 
We call the quantity $$\sum_{i<j}|t_i-t_j|$$
the {\it total width} of $L$ with respect to $\ttt$.   
Let  $\left\{x_{d_1,1}=x_{d_1,c_1},\ldots,x_{d_v,c_v}=x_{k,c_v}\right\}$ be the lower-outside corners of $L$. We say that a pair $(L', \ttt')$ with $L'\subset X$ a ladder and $\ttt'\in \NN^d$  is the {\it chamfer} of $(L,\ttt)$ at $x_{d_j,c_j}$, or that $(L', \ttt')$ is obtained from $(L,\ttt))$ by {\it chamfering} $x_{d_j,c_j}$, if $\ttt'=(t_1,\ldots, t_{j-1}, t_j-1,t_{j+1},\ldots, t_v)$, and $L'$ has the same upper-outside corners of $L$ and lower-outside corners $$\left\{x_{d_1,1}=x_{d_1,c_1},\ldots,x_{d_{j-1},c_{j-1}},\, x_{d_j-1,c_j+1},\, x_{d_{j+1},c_{j+1}},\ldots,x_{d_v,c_v}=x_{k,c_v}\right\}$$ for some  %a lower-outside corner 
$x_{d_j,c_j}$ such that $d_j<d_{j-1}$ and $c_j>c_{j-1}$. 
\end{definition}

%The following proposition with its corollaries are  the main results of this section.

\begin{proposition} \label{prop_chamfer}
Assume \autoref{setup_chamfer}, and let $L$ be a non-degenerate ladder with respect to $\ttt$. Then
\begin{enumerate}[{\rm (1)}]
	\item if $L'$ is the chamfer of $(L,\ttt)$ at $x_{d_j,c_j}$, then  $L'$  is non-degenerate.
	\item $L$ is the chamfer of a non-degenerate ladder with strictly smaller total width. 
\end{enumerate}
\end{proposition}
\begin{proof} We note that condition (3) of \autoref{running assumptions ladder} holds if and only if $d_{j+1}-d_j>t_{j+1}-t_j$ and $c_{j+1}-c_j>t_{j}-t_{j+1}$, and one can readily check that it holds for $L'$. The fact that $L'$ satisfies the other conditions of \autoref{running assumptions ladder} is an easy check. For the second claim, let $j$ be such that $t_j=\min\{t_1,\ldots, t_v\}$. Consider the ladder $L''$ with same upper-outside corners as $L$ and lower outside corners $$\left\{x_{d_1,1}=x_{d_1,c_1},\ldots,x_{d_{j-1},c_{j-1}},\, x_{d_j+1,c_j-1},\, x_{d_{j+1},c_{j+1}},\ldots,x_{d_v,c_v}=x_{k,c_v}\right\}.$$
We note that $L''$ is non-dengenerate with respect to $\ttt''=(t_1,\ldots, t_{j-1}, t_j+1,t_{j+1},\ldots, t_v)$, $(L,\ttt)$ is the chamfer of $(L'',\ttt'')$ at $x_{d_j+1,c_j-1}$, and the total width of $L''$ is strictly smaller than $L$.
\end{proof}

\begin{corollary} \label{coroll_chamfer}
Assume \autoref{setup_chamfer}. If $L$ is  non-degenerate with respect to $\ttt$, then there exists a ladder $L''$ and $t,u\in \NN$ such that $L''$ is non-degenerate with respect to $(t,\ldots, t)\in \NN^u$ and $L$ is obtained from $L''$ by a finite sequence of chamferings. 
\end{corollary}
\begin{proof}
We proceed by descending induction on the total width of the ladder, making a repeated use of \autoref{prop_chamfer} (2).
\end{proof}

For an ideal $I$ in a commutative ring with identity $A$, we denote the {\it symbolic Rees algebra} of $I$ by $\R^s(I)=\oplus_{n\gs 0}I^{(n)}T^n\subseteq A[T]$, where $T$ is a variable. Sometimes we write  $\R^s_A(I):=\R^s(I)$ if the ambient ring needs to be specified. We denote by $\gr^s(I)=\oplus_{n\gs 0}I^{(n)}/I^{(n+1)}$ the {\it symbolic associated graded algebra} of $I$. 

\begin{theorem} \label{ThmNoetherian}
Assume \autoref{setup_chamfer}, and let $L$ be a non-degenerate ladder with respect to $\ttt$. Then $\R^s(I_\ttt(L))$ and $\gr^s(I_\ttt(L))$ are Noetherian rings.
\end{theorem}
\begin{proof}
It suffices to show the statement for $\R^s(I_\ttt(L))$. By Corollary \ref{coroll_chamfer} there exists a ladder $L''$ which is non-degenerate with respect to $(t,\ldots,t) \in \NN^u$, and $L$ is obtained from $L''$ by a finite sequence of chamferings. Then by Gorla's work on mixed ladder determinantal ideals \cite[Lemma 1.19]{GorlaMixed} there are indeterminates $\{y_1,\ldots,y_d\}$ of $R''=\kk[L'']$ and $c \leq d$ such that 
\[
I_\ttt(L)[y_1,\ldots,y_d,y_1^{-1},\ldots y_c^{-1}] \cong I_t(L'')[y_1^{-1},\ldots y_c^{-1}].
\]
From this isomorphism we conclude that for all $n \geq 1$
\[
I_\ttt(L)^{(n)}[y_1,\ldots,y_d,y_1^{-1},\ldots y_c^{-1}] \cong I_t(L'')^{(n)}[y_1^{-1},\ldots y_c^{-1}].
\]
Since $\R^s_{R''}(I_t(L''))$ is Noetherian \cite{brunsconca98}, so is $\R^s_R(I_\ttt(L))[y_1,\ldots,y_d,y_1^{-1},\ldots y_c^{-1}]$. As $\R^s_R(I_\ttt(L))$ is a direct summand of the latter, we conclude that it is itself Noetherian.
\end{proof}

\begin{remark} Unnmixed ladder determinantal ideals define $F$-rational singularities in characteristic $p\gg0$, and rational singularities in characteristic zero, by work of Conca and Herzog \cite{concaherzog1997}. It follows from \autoref{coroll_chamfer} that this is also true for unmixed ladder determinantal ideals since being $F$-rational or having rational singularities localize, and because if $S$ is a ring and $x$ is an indeterminate one has that $S$ is $F$-rational (resp. has rational singularities) if and only if $S[x]$ is $F$-rational (resp. has rational singularities).
\end{remark}

%%%%%%%%%%%%%%%%%%%%%%%%%%%%%%%%%%%%%%%%%%%%%%%%%%%%%%%%%
\section{Mixed ladder determinantal ideals}\label{MLDSection}
%%%%%%%%%%%%%%%%%%%%%%%%%%%%%%%%%%%%%%%%%%%%%%%%%%%%%%%%%


In this section we prove the main results of this paper regarding the $F$-singularities of ladder determinantal rings and their blowup algebras. Throughout we use the following setup:

%To prove the next theorem, it is useful to introduce some notation:

\begin{setup} \label{setup_mainsec_ladder} 
We adopt \autoref{setup_chamfer}, and when the base field $\kk$ has characteristic $p>0$ we further assume that it is $F$-finite. For  $r=2,\ldots ,k+\ell$ we denote by $\mathcal{D}_r$ the antidiagonal at level $r$ of $X$, namely, 
\[\mathcal{D}_r=\{x_{i,j}\in X\mid i+j=r\}\subseteq X.\]
% Given  $\ttt=(t_1,\ldots,t_v) \in \ZZ_{>0}^v$, 
We define $L^\circ_{\ttt} \subseteq L$ to be the subladder with the same upper-outside corners as $L$, and with lower-outside corners  $\{x_{d_j-t_j+1,c_j+t_j-1} \mid 1 \leq j \leq v\}$. Notice that $L^\circ_{\ttt} =\bigcup_{i=1}^v (L_j)^\circ_{t_j}$ (see \autoref{eq_Lj}).
\end{setup}
%By a result of Gorla, we have that $\Ht(I_t(L)) = |L^\circ|$ \cite[Theorem 1.15]{GorlaMixed}. 

The following is the first main result of this section. 

\begin{theorem}\label{ThmMixLadderSFP}
Assume \autoref{setup_mainsec_ladder} with   $\mathbb{k}$ of  characteristic $p>0$. %Let $X$ be a $k \times l$ matrix of indeterminates and $L\subseteq  X$ a ladder. For $\ttt=(t_1,\ldots,t_v) \in \ZZ_{>0}^v$, we let $I_{\ttt}(L) \subseteq R=\mathbb{K}[L]$ be the corresponding mixed ladder determinantal ideal. 
Then $I_{\ttt}(L)$ is symbolic $F$-split.
\end{theorem}
\begin{proof}
Let $h = \Ht\left(I_{\ttt}(L)\right)$. It suffices to find an element $f \in I_{\ttt}(L)^{(h)}$ and a monomial order $\prec$ such that $\IN_\prec(f)$ is squarefree \cite[Lemma 6.2(1)]{DSMNB}. 
Let 
\begin{equation}\label{eq_inv1}
\mathcal{A} = \{2 \leq r \leq k+\ell \mid \mathcal{D}_r\cap L^\circ_\ttt\neq\emptyset\},
\end{equation}
and for every $r \in \mathcal{A}$  set 
\begin{equation}\label{eq_inv2}
w_r=\max\{w\mid x_{w,r-w}\in \mathcal{D}_r\cap L^\circ_\ttt\}
\qquad
\text{and}
\qquad
p_r = \max\{j \mid x_{w_r, r-w_r}\in  (L_j)^\circ_{t_j} \}.
\end{equation}
Furthermore,  set 
\begin{equation}\label{eq_inv3}
a_r=\min\left\{i\mid x_{i,r-i}\in L_{p_r}\cap \mathcal{D}_r\right\} \qquad \text{and}\qquad b_r=\max\left\{i\mid x_{i,r-i}\in L_{p_r}\cap \mathcal{D}_r\right\}.
\end{equation}
Let $\gamma_r= b_r-a_r+1$ and observe that the $\gamma_r \times \gamma_r$ matrix $Y_r=X^{[a_r,b_r]}_{[r-b_r,r-a_r]}$ is contained in $L_{p_r}$ (see \autoref{notation_affine_subsec}). 
  We note that %$\det(Y_r) \in I_{\ttt}(L)^{\left(c_r-t_{p_r}+1\right)}$. Indeed, since $x_{w_r,r-w_r} \in \Lambda_{p_r}$ we have that $Y_r \subseteq L_{p_r}$. It follows that 
\begin{equation}\label{eq_detY}
\det(Y_r) \in I_{\gamma_r}(Y_r) \subseteq I_{t_{p_r}}(Y_r)^{\left(\gamma_r-t_{p_r}+1\right)} \subseteq I_{t_{p_r}}(L_{p_r})^{\left(\gamma_r-t_{p_r}+1\right)} \subseteq I_{\ttt}(L) ^{\left(\gamma_r-t_{p_r}+1\right)},	
\end{equation}
where the first inclusion follows from \cite[Proposition 10.2]{BookDet}.


Let $\mathcal{B}=\{r \in \mathcal{A}\mid \gamma_r-t_{p_r}+1 \geq 0\}$ and set 
\begin{equation}\label{eq_f}
f=\prod_{r \in \mathcal{B}} \det(Y_r)
\end{equation}
 We claim that $f \in I_{\ttt}(L)^{(h)}$. Observe that for every $r \in \mathcal{B}$, the integer $\max\{\gamma_r-t_{p_r}+1,0\}$ counts  the number of elements on the main antidiagonal of $Y_r$ that belong to $L^\circ_\ttt$. 
Moreover, %By our choice of the matrices $Y_r$, 
every entry of $L^\circ_\ttt$  belongs to the main antidiagonal of  $Y_r$ for exactly one $r \in \mathcal{A}$. Thus, by  \cite[Theorem 1.15]{GorlaMixed} we  have 
\begin{equation}\label{eq_sum_height}
	\sum_{r \in \mathcal{B}}(\gamma_r-t_{p_r} + 1) = |L^\circ_\ttt| = h.
\end{equation}
 Therefore, \autoref{eq_detY} implies
\[
f \in \prod_{r \in \mathcal{B}} I_{\ttt}(L)^{\left(\gamma_r-t_{p_r}+1\right)} \subseteq I_{\ttt}(L)^{(h)},
\]
whence the claim follows. Finally, if $\prec$ is any antidiagonal term order, then $\IN_\prec(f)$ is squarefree, completing the proof.
\end{proof}

As a consequence of \autoref{ThmMixLadderSFP} we answer a question of Conca and Herzog on $F$-purity of algebras defined by ladder determinantal ideals \cite[p. 122]{concaherzog1997}. Much more generally, we obtain that symbolic blowup algebras of mixed ladder determinantal ideals are $F$-pure algebras. 

\begin{corollary}\label{blowups_f_pure}
Assume \autoref{setup_mainsec_ladder} with   $\mathbb{k}$ of  characteristic $p>0$.
Then the blowup algebras $\R^s\left(I_{\ttt}(L)\right)$ and $\gr^s\left(I_{\ttt}(L)\right)$ are $F$-pure. In particular, $R/I_\ttt(L)$ is $F$-pure.
\end{corollary}
\begin{proof}
The first claim follows from  \autoref{ThmMixLadderSFP} and properties of symbolic $F$-split ideals \cite[Theorem 4.7]{DSMNB}. The second follows the fact that $R/I_\ttt(L)$ is a direct summand of $\gr^s(I_\ttt(L))$.
\end{proof}

Let $S=\kk[x_1,\ldots, x_d]$ be a standard graded polynomial. For a finitely generated $\ZZ$-graded $S$-module $M$ we  denote by $\reg(M)$ its   {\it Castelnuovo-Mumford regularity}.

\begin{corollary}\label{cor_RegDEpth}
Assume \autoref{setup_mainsec_ladder} with   $\mathbb{k}$ of  characteristic $p>0$.  Then,
$\lim\limits_{n\to\infty}\frac{\reg\left(R/I_\ttt(L)^{(n)}\right)}{n}$ exists and the sequence
$\{\Depth(R/I_\ttt(L)^{(n)})\}_{n \in \NN}$ is constant for $n\gg0$, with value equal to $\min\{\depth(R/I_\ttt(L)^{(n)})\}_{n \in \NN}$.
\end{corollary}
\begin{proof}
The first  statement follows from  \autoref{ThmMixLadderSFP} and properties of symbolic $F$-split ideals \cite[Theorem 4.10]{DSMNB}, because the symbolic Rees algebra of $I_\ttt(L)$ is Noetherian \autoref{ThmNoetherian}. 
\end{proof}

%For the convenience of the reader, we record the following statement which  combines results from previous work \cite{DSMNB,KoleyVarbaro}.
%
%\begin{proposition}\label{symbolic<-x>initial}
%Let $I\subseteq S=\mathbb{k}[x_1,\ldots ,x_d]$ be a homogeneous prime ideal of height $h$ and $\prec$ a term order on $S$.  If there exists $f\in I^{(h)}$ with $\IN_\prec(f)$ squarefree, then $\IN_\prec(I)$ is a squarefree monomial ideal. Furthermore, if $\init_\prec\left(\cR^s(I)\right)$ is generated as an $S$-algebra by forms of degrees $b_1,\ldots ,b_m$ and $\mathbb{k}$ has characteristic $p>\lcm(b_1,\ldots ,b_m)$, then $\IN_\prec(I^{(n)})=(\IN_\prec(I))^{(n)}$ for all $n\in\NN$.
%\end{proposition}
%\begin{proof}
%The monomial $\IN_\prec(I)$ is squarefree \cite[Theorem 3.13]{KoleyVarbaro}. 
%It follows that the graded algebra   $\bigoplus_{n\in\NN}\IN_\prec(I^{n})/\IN_\prec(I^{n+1})$ is reduced  \cite[Proposition 7.5]{DSMNB}. 
%Finally, we conclude that $\IN_\prec(I^{(n)})=(\IN_\prec(I))^{(n)}$ for all $n\in\NN$
%\cite[Corollary 7.10]{DSMNB}.
%\end{proof}

We also show that symbolic powers commute with taking inital ideals for unmixed ladder determinantal ideals.
\begin{corollary}\label{cor_SAGBI}
Assume \autoref{setup_mainsec_ladder} with   $\mathbb{k}$ a perfect field of  characteristic $p>0$. For $t \in \ZZ_{>0}$, we let $I_t(L) \subseteq R=\mathbb{k}[L]$ be the corresponding unmixed ladder determinantal ideal.  If $\prec$ denote an antidiagonal order and $p>(\min\{k,\ell\}-t+1)!$, then
$$
\IN_\prec\left(I_t(L)^{(n)}\right)=\left(\IN_\prec \left(I_t(L)\right)\right)^{(n)} \,\,\, \text{ for every }n\in\NN.
$$
\end{corollary}
\begin{proof}
We have that
$\init_\prec\left(\cR^s\left(I_t(L)\right)\right)$  is generated by forms of degree of degree at most $\min\{k,\ell\}-t+1$  \cite[Theorem 4.1]{brunsconca98}, so the last statement follows from  \autoref{symbolic-initial} and the Remark following it.
\end{proof}

%\begin{remark}
%It is known  that taking  symbolic powers and initial ideals of determinantal ideals are commuting operations  in any characteristic \cite[Proposition 7.2]{BCsymbinit}. From  \autoref{cor_SAGBI}  it follows that this fact holds more generally for unmixed ladder determinantal ideals. We note that our proof is different in nature as the one previously known \cite{BCsymbinit}, which    involves a careful combinatorial analysis.
%\end{remark}

The following technical lemma is needed for the proof of Theorem \ref{stronglyF}.

\begin{lemma}\label{Lemma sfr} Let $B$ be an $F$-finite ring of characteristic $p>0$, $J \subseteq B$ an ideal, and $A=B[x_1,\ldots,x_d]$ a polynomial extension. Let $H=(x_1,\ldots,x_d)$ and  $I=JA+H$. Then $\cR^s_A(I)$ is strongly $F$-regular if and only if $\cR^s_B(J)$ is strongly $F$-regular.
\end{lemma}
\begin{proof}
Since $H^{(j)}=H^j$ for all $j \in \ZZ_{>0}$,   we have that $I^{(j)} = \sum_{i=0}^j (JA)^{(i)}H^{j-i}$  \cite[Theorem 3.4]{BinomialSumSP}. Thus
%\[
%\cR^s_R(I) =\bigoplus_{j \in \NN} I^{(j)}T^j =  \bigoplus_{j \in \NN} \left(\sum_{i=0}^j (JR)^{(i)}T^iH^{j-i}T^{j-i}\right) \cong \cR^s_R(JR)[y_1T,\ldots,y_mT].
%\]
\[
\cR^s_A(I) = \cR^s_A(JA)[x_1T,\ldots,x_dT]=\cR^s_B(J)[x_1,\ldots,x_d,x_1T,\ldots,x_dT],
\]
where the second equality holds as   $(JA)^{(i)} = J^{(i)}A = J^{(i)}[x_1,\ldots,x_d]$. Thus $\cR^s_A(I)$ is isomorphic to $\cR^s_B(J)[x_1Z,\ldots,x_dZ,x_1T,\ldots,x_dT]$, with $Z$ another variable, and so isomorphic to the Segre product  $\cR^s_B(J)[x_1,\ldots,x_d]\,
{\mathlarger{\mathlarger\#}}\,\cR^s_B(J)[Z,T]$. % and so isomorphic to a direct summand of the latter algebra

Now, assume  $\cR^s_A(I)$ is strongly $F$-regular. Since $\cR^s_B(J)$ is a direct summand of $\cR^s_A(I)$, it is strongly $F$-regular  \cite[Theorem 3.1(e)]{HHStrongFreg}. Conversely, if $\cR^s_B(J)$ is strongly $F$-regular then so is the polynomial extension $\cR^s_B(J)[x_1,\ldots,x_d,Z,T]$. Since $\cR^s_A(I)$ is isomorphic to %the Segree product  $\cR^s_B(J)[y_1,\ldots,y_m]\,
%{\mathlarger{\mathlarger\#}}\,\cR^s_B(J)[Y,Z]$, and so isomorphic to 
a direct summand of the latter, it is  strongly $F$-regular as well, finishing the proof. 
\end{proof}

The following is one of main result of this paper. We show that the symbolic Rees algebra of any mixed ladder determinantal ideal is strongly $F$-regular. 

\begin{theorem}\label{stronglyF}
Assume \autoref{setup_mainsec_ladder} with   $\mathbb{k}$ of characteristic $p>0$. Then  $\R^s\left(I_{\ttt}(L)\right)$ is strongly $F$-regular.
\end{theorem}
\begin{proof}
%If we prove the theorem for the base field $\overline{\mathbb{K}}$ then the theorem will also follow for $\mathbb{K}$ \cite[Lemma 3.17]{Hashimoto}, therefore we may assume without loss of generality that $\mathbb{K}$ is perfect. 

We proceed by induction on $d=\dim(R)$, i.e., on the number of entries of. The base case $d=1$ is trivial.

It $t_v=1$ then we can write $I_{\ttt}(L)=I_{\ttt'}(L') + I_1(L_v)$ where $L'$ is a proper subladder of $L$ and $\ttt' = (t_1,\ldots,t_{v-1})$. %We have that $\cR^s(I)$ is a polynomial extension of $\cR^s(I_{\ttt'}(L'))$ \cite[Theorem 3.4]{BinomialSumSP}, and if the latter is strongly F-regular then so is the former. 
Let $R' = \mathbb{k}[L']$. Since $\dim(R')<\dim(R)$ by induction hypothesis  we have that $\cR^s_{R'}\left(I_{\ttt'}(L')\right)$ is strongly $F$-regular. It follows that $\cR^s_R\left(I_{\ttt}(L)\right)$ is strongly $F$-regular by \autoref{Lemma sfr}.

Now assume $t_v>1$ and  consider the smallest $\alpha \in \{1,\ldots,v\}$ for which $(k,c_\alpha)$ is a lower-outside corner of $L$. We observe that such a $\alpha$ exists since $(k,c_v)$ is always a lower-outside corner by our  assumptions (see Notation \ref{running assumptions ladder}). %We use the notation introduced before \autoref{ThmMixLadderSFP}.
Let $\mathcal{A},p_r, a_r, b_r$ be as in \autoref{eq_inv1}, \autoref{eq_inv2}, and \autoref{eq_inv3}. For every $r\in\mathcal{A}$ set $\gamma_r= b_r-a_r+1$ and  $Y_r=X^{[a_r,b_r]}_{[r-b_r,r-a_r]}$. 
 By \autoref{eq_detY} we have $\det(Y_r) \in I_\ttt(L)^{\left(\gamma_r-t_{p_r}+1\right)}$.  
 Let $\mathcal{B}=\{r \in \mathcal{A}\mid \gamma_r-t_{p_r}+1 \geq 0\}$   and note that $\beta:=k+\gamma_\alpha \in \mathcal{B}$, $b_\beta = k$,  and $p_\beta = v$. 
 Set $Y=X^{[a_{\beta},k-1]}_{[c_\alpha+1,\beta-a_{\beta}]}$ and observe that 
$ \det(Y) \in I_{\gamma_\beta-1}(Y) \subseteq I_{t_v}(Y)^{(\gamma_\beta-t_v)}\subseteq  I_\ttt(L)^{(\gamma_\beta-t_v)}$  \cite[Proposition 10.2]{BookDet}.
Set $$
\mathcal{B}\,' = \mathcal{B} \smallsetminus \{\beta\}
\qquad 
\text{and} 
\qquad
g=\det(Y)  \prod_{r \in \mathcal{B}\,'} \det(Y_r).$$ We claim that $g \in I_\ttt(L)^{(h-1)}$, where $h=\Ht\left(I_\ttt(L)\right)$. 
For every $r \in \mathcal{B}$, the integer $\max\{\gamma_r-t_{p_r}+1,0\}$ counts  the number of elements on the main antidiagonal of $Y_r$ that belong to $L^\circ_\ttt$. By \autoref{eq_sum_height} we  have
 $$(\gamma_{\beta}-t_v)+\sum_{r \in \mathcal{B}\,'}(\gamma_r-t_{p_r} + 1) =\sum_{r \in \mathcal{B}}(\gamma_r-t_{p_r} + 1)=  h-1.$$ 
 Therefore $g \in I_\ttt(L)^{(\gamma_{\beta}-t_v)}  \prod_{r \in \mathcal{B}\,'} I_\ttt(L)^{(\gamma_r-t_{p_r}+1)} \subseteq I_\ttt(L)^{(h-1)}$, which proves the claim. 
 

 
  Now let $\prec$ be any antidiagonal term order and note that  $\IN_\prec(g)$ is a squarefree monomial not divisible by $x_{k,c_\alpha}$. Moreover,  we have $g^{p-1}\in \left(I_{\ttt}(L)^{(n)}\right)^{[p]}:I_{\ttt}(L)^{(np)}$ for every $n\in \NN$ \cite[Lemma 2.6]{GrifoHuneke}. Thus there exists an $\cR^s_R\left(I_{\ttt}(L)\right)$-homomorphism
  $\Psi:\cR^s_R\left(I_{\ttt}(L)\right)^{1/p}\to\cR^s_R\left(I_{\ttt}(L)\right)$ such that 
 $\Psi\left(x_{k,c_\alpha}^{1/p}\right)=1$ (cf. proof of \cite[Theorem 6.7]{DSMNB}). 

  
  
Since $\R^s(I_\ttt(L))$ is Noetherian by \autoref{ThmNoetherian}, and $\kk$ is $F$-finite, we have that $\R^s(I_\ttt(L))$ is $F$-finite. By  \cite[Theorem 3.3]{HHStrongFreg} to finish the proof it suffices to show that $\cR^s_R(I_{\ttt}(L)^{(n)})_{x_{k,c_\alpha}}$ is strongly $F$-regular. Let $L''$ be the ladder obtained from $L$ by deleting the entries $$E=\{x_{d_{\alpha-1}+1,c_\alpha},x_{d_{\alpha-1}+2,c_\alpha},\ldots,x_{k,c_\alpha},x_{k,c_\alpha+1},\ldots,x_{k,c_{\alpha+1}-1}\}$$ from $L$. Set $R''=\mathbb{k}[L'']$ and $S=R''[E]_{x_{k,c_\alpha}}$.  There is an isomorphism 
  $$\theta:R_{x_{k,c_\alpha}} \to S 
\quad \text{such that} 
\quad 
\theta\left(I_\ttt(L)_{x_{k,c_\alpha}}\right) = I_{\ttt''}(L'')S,
  $$ 
  where $\ttt''=(t_1,\ldots,t_{\alpha -1},t_\alpha -1,t_{\alpha +1}, t_{v-1},t_v)$ and $I_{\ttt''}(L'')$ is computed in $R''$  \cite[Lemma 1.19]{GorlaMixed}. One can check that  $\theta\left(I_\ttt(L)^{(n)}_{x_{k,c_\alpha}}\right) = I_{\ttt''}\left(L''\right)^{(n)}S$ for all $n \geq 1$ (see \cite[Lemma 10.1]{BookDet} for analogue procedure). By induction hypothesis  $\cR^s_{R''}(I_{\ttt''}(L''))$ is strongly $F$-regular, and therefore $\cR^s_{R''}(I_{\ttt'}(L'')) \otimes_{R''} S$ is strongly $F$-regular. Therefore we conclude via the  isomorphism $\theta$ that $\cR^s_R(I_{\ttt}(L)^{(n)}) \otimes_R R_{x_{k,c_\alpha}} \cong \cR^s_R(I_{\ttt}(L)^{(n)})_{x_{k,c_\alpha}}$ is strongly $F$-regular, finishing the proof.
\end{proof}




%%%%%%%%%
The rest of the this section is devoted to proving the following theorem which states  that unmixed determinantal ideals of two-sided ladders are Knutson ideals. Following \autoref{notation_affine_subsec}, for a ladder $L\subset X$ we denote by  $L_{[a,b]}=X_{[a,b]} \cap L$ the  subladder corresponding to the  consecutive columns from $a$ to $b$. Likewise,  $L^{[c,d]}=X^{[c,d]} \cap L$ denotes the subladder corresponding to the  consecutive rows from $c$ to $d$. 

\begin{theorem}\label{ThmLadderKI}
Assume \autoref{setup_mainsec_ladder}. For $t \in \ZZ_{>0}$  we let $I_t(L) \subseteq R=\mathbb{k}[L]$ be the corresponding unmixed ladder determinantal ideal. Let 
 $f$ be as in \autoref{eq_f} with $t_1=\cdots =t_v=t$. Then $$I_t\left(L_{[a,b]}\right) \in \mathcal{C}_f \qquad 
\text{ for every } \qquad 1\leq a < b\leq \ell \quad \text{ such that } \quad I_t\left(L_{[a,b]}\right)\neq 0,\qquad \text{and}$$   $$I_t\left(L^{[c,d]}\right) \in \mathcal{C}_f \qquad  
\text{ for every } \qquad 1\leq c < d\leq k \quad \text{ such that } \quad I_t\left(L^{[c,d]}\right)\neq 0.$$
In particular, then $ I_t(L)\in \mathcal{C}_f$ whenever $I_t(L)$ is not zero.
\end{theorem}


 A first step towards the proof of  \autoref{ThmLadderKI} is the following lemma where we show that the ideal generated by the $t$-minors of a subladder corresponding to $t$ consecutive columns or rows is in $\mathcal{C}_f$.
\begin{lemma}\label{lemmaLadderKI}
Under the assumptions and notations as  in  \autoref{ThmLadderKI} we have
$$
 I_t\left(L_{[j,j+t-1]}\right) \in \mathcal{C}_f \qquad  \text{ for every } \,\,\,\, 1\ls j\ls  \ell-t+1 \quad \text{ such that } \quad I_t\left(L_{[j,j+t-1]}\right)\neq 0,\qquad \text{and}
 $$
 $$
 I_t\left(L^{[i,i+t-1]}\right) \in \mathcal{C}_f \qquad\text{ for every } \qquad 1\ls i\ls k-t+1 \quad \text{ such that } \quad I_t\left(L^{[i,i+t-1]}\right)\neq 0.
 $$
\end{lemma}

\begin{proof}
Fix $1\ls j\ls \ell-t+1$ such that $I_t\left(L_{[j,j+t-1]}\right)\neq 0$. Set $L'=L_{[j,j+t-1]}$ and  $\mathcal{E}=\{2 \leq r \leq k+\ell \mid \left|\mathcal{D}_{r}\cap L'\right|=t \}$. %. By construction of $L'$, $\mathcal{A}^{\prime}=\{2 \leq r^{\prime} \leq k+l \mid |D_r\cap \widetilde{L}|=t \}$. 
Notice that $\mathcal{E} \neq 0$, and for every $r \in \mathcal{E}$ there is only  one element on the main antidiagonal of $Y_{r}$ that belongs to $(L')^\circ_t$ (see sentence after \autoref{eq_inv3}). Then by  \cite[Theorem 1.15]{GorlaMixed} it follows that $\Ht \left(I_t(L')\right)=|(L)^\circ_t| = | \mathcal{E}|$. Moreover
\begin{equation*}\label{eq_dets_in}
	H:=\left( \left\{\det(Y_{r}) \mid r \in \mathcal{E}\right\}\right)\subseteq I_t(L')  . 
\end{equation*}
%Define $H$ to be the ideal on the right hand side of \autoref{eq_dets_in} and observe %. Let $\prec$ be a an antidiagonal term order and observe that the initial ideal of $H$ is given by the product of the elements in some of the main antidiagonals of $L'$. Since these monomials are coprime, this ideal 
Notice $H$ is a complete intersection with $\Ht (H)=|\mathcal{E}|$ and $H\in \mathcal{C}_f$ as it is generated by some of the factors of $f$. Thus $I_t(L')$ is a minimal prime of  of $H$ and so $I_t(L')\in \mathcal{C}_f$. 

Likewise, one can prove that $I_t\left(L^{[i,i+t-1]}\right)$ is either zero or in $ \mathcal{C}_f$ for every $i$.
\end{proof}

With \autoref{lemmaLadderKI} in hand we can  prove  \autoref{ThmLadderKI}.

\begin{proof}[Proof of \autoref{ThmLadderKI}]
%Let $t$ be such that $\Ht \left(I_t(L)\right)>0$. % We want to prove that $I_t(L) \in \mathcal{C}_f$. 
We only show $I_t\left(L_{[a,b]}\right) \in \mathcal{C}_f$ whenever $I_t\left(L_{[a,b]}\right) \neq 0$, as the proof for $I_t\left(L^{[c,d]}\right) $ is identical.


We proceed by induction on $\delta\in \NN$ to show that $I_t\left(L_{[j,j+t-1+\delta]}\right)$ is either 0 or in $\mathcal{C}_f$  for every $1\ls j\ls \ell-t+1-\delta$, the base case $\delta=0$ being 
\autoref{lemmaLadderKI}. %Notice that since $I_t\left(L_{[a,b]}\right) \neq 0$, then there always exists an index $j$ such that $0 \neq I_t\left(L_{[j,j+t-1]}\right)\in \mathcal{C}_f$.

Fix $1\ls j\ls \ell-t-\delta$. By  induction hypothesis  the ideals $I_t\left(L_{[j,j+t-1+\delta]}\right)$ and $I_t\left(L_{[j+1,j+t+\delta]}\right)$ are either zero or in $\mathcal{C}_f$. Therefore so is their  sum. Assume that the sum is not zero.  

We claim that 
\begin{equation}\label{eq_first_int}
I_t\left(L_{[j,j+t-1+\delta]}\right)+I_t\left(L_{[j+1,j+t+\delta]}\right)=I_t\left(L_{[j,j+t+\delta]}\right) \cap I_{t-1}\left(L_{[j+1,j+t-1+\delta]}\right),
\end{equation}
if  $t>1$, %and both the summands are not zero, 
 and  
%$I_{t-1}\left(L_{[j+1,j+t-1+\delta]}\right)\neq (0)$, and
\begin{equation}\label{eq_second_int}
	I_t\left(L_{[j,j+t-1+\delta]}\right)+I_t\left(L_{[j+1,j+t+\delta]}\right)=I_t\left(L_{[j,j+t+\delta]}\right),
\end{equation}
othwewise.

Since \autoref{eq_second_int} is clear, we focus on proving \autoref{eq_first_int}. So, we may assume  $t>1$.  % and   $I_{t-1}\left(L_{[j+1,j+t-1+\delta]}\right)\neq (0).$ 
To simplify the notation we set 
$$L':=L_{[j,j+t-1+\delta]}\qquad \text{and} \qquad L'':=L_{[j+1,j+t+\delta]}.$$ 
% Thus, we want to prove that $$I_t(L_1)+I_t(L_2)=I_t(L_1 \cup L_2) \cap I_{t-1}(L_1 \cap L_2).$$
 From the proof of \cite[Theorem~2.1]{seccia2022knutson} it follows that $$I_t\left(X_{[j,j+t-1+\delta]}\right)+I_t\left(X_{[j+1,j+t+\delta]}\right)=I_t\left(X_{[j,j+t+\delta]}\right) \cap I_{t-1}\left(X_{[j+1,j+t-1+\delta]}\right).$$
 Thus after contracting these ideals to $\mathbb{k}[L' \cup L'']$ we obtain
$$
  \left( I_t\left(X_{[j,i+t-1+\delta]}\right)+I_t\left(X_{[j+1,j+t+\delta]}\right)\right) \cap \mathbb{k}[L' \cup L'']=I_t(L') + I_{t}(L''),\qquad\text{and}$$
  $$
  \left(I_t\left(X_{[j,j+t+\delta]}\right) \cap I_{t-1}\left(X_{[j+1,j+t-1+\delta]}\right)\right)\cap \mathbb{k}[L' \cup L'']= I_t(L' \cup L'') \cap I_{t-1}(L' \cap L''),
$$
where the first equality follows from \cite[Lemma 2.2]{seccia2022knutson}, \cite[Corollary 2]{knutson2009frobenius}, and  \cite[Lemma 1.8]{GorlaMixed}, whereas the second one is clear.  
 %$I_t(X_{[i,i+t-1]})$ and $I_t(X_{[i+1,i+t]})$ are Knutson ideals.
 This proves the claim. 
 
 Hence the ideal
%
%
%We point out that the intersection $I_t(L' \cup L'') \cap I_{t-1}(L' \cap L')$ could be redundant in some cases. If this happens, we have $I_t(L' \cup L'')=I_t(L')+I_t(L'') \in \mathcal{C}_f$. Otherwise, 
$I_t(L_{[j,j+t+\delta]})=I_t(L' \cup L'')$ is a minimal prime of $I_t(L')+I_t(L'') \in \mathcal{C}_f$. We conclude  that $I_t(L_{[j,j+t+\delta]})$ is either 0 or in $\mathcal{C}_f$ for every $1\ls j\ls  \ell-t$, finishing the proof.  
\end{proof}


%%%%%%%%%%%%%%%%%%%%%%%%%%%%%%%%%%%%%%%%%%%%%%%%%%%%%%%%%
\section{Knutson  Schubert determinantal and poset ideals}\label{SectionSchubertMSV}
%%%%%%%%%%%%%%%%%%%%%%%%%%%%%%%%%%%%%%%%%%%%%%%%%%%%%%%%%

In this section, we show that Schubert determinantal ideals and poset determinantal ideals are Knutson ideals. We remark that the statement for  Schubert determinantal ideals  already appeared in Knutson's work, but our proof uses different methods \cite{knutson2009frobenius}. %We begin by recalling the following defintions. 

Throughout this section we follow the notations introduced in \autoref{sub_affine_schub}. In particular, $X=(x_{i,j})$ is a generic matrix of size $k \times \ell$. % with $k\ls \ell$
We also assume the following setup.

\begin{setup}\label{setup_Knutson_section}
	We adopt  \autoref{notation_affine_subsec}. % and   denote by 
	%	Let $X=(x_{i,j})$ be a generic matrix of size $k \times \ell$. 
	Let $R=\kk[X]$ be the polynomial ring in the variables in $X$, where $\kk$ is a an arbitrary field. Moreover, we equip $R$ with an antidiogonal term order.
\end{setup}




We continue with the definition of the following polynomial which plays an essential role in the results in this section. This polynomial is inspired by the results by the fourth-named author \cite{seccia2022knutson}.
\begin{definition}\label{def_f(X)}
	Assume  \autoref{setup_Knutson_section}. We define $f(X) \in \kk[X]$ as: %to be the polynomial given by the product of all the minors of $X$ corresponding to its antidiagonals, that is 
	$$f(X)= \prod_{j=0}^{k-2} \left[\det \left(X_{[1,j+1]}^{[1,j+1]}\right) 
	\det \left(X^{[k-j,k]}_{[\ell-j,\ell]}\right) \right] 
	\prod_{j=1}^{\ell-k+1} \left(\det \left(X_{[j,k+j-1]}\right) \right).$$ 
	We note that: 
	$$\init (f(X))= \prod \limits_{1\ls i\ls k,1\ls j\ls \ell} x_{i,j}.$$
	%Therefore, we can construct the Knutson family of ideals associated with this $f$. This choice of $f$ allows us to use some known results from \cite. In particular, we will use the fact that given a generic matrix, all determinantal ideals on adjacent columns (or adjacent rows) are Knutson ideals for this (standard) choice of $f$ (see \cite[Theorem 2.1] for more details).
\end{definition}

The following theorem shows that determinantal ideals of every northwest and southeast submatrix of $X$ is Knutson associated to $f(X)$. This allows us to show Schubert determinantal and poset ideals are Knutson

\begin{theorem}\label{p:msk}
	Assume \autoref{setup_Knutson_section} and  let $f:=f(X) $ be as in \autoref{def_f(X)}.  Then for every $1\ls r \ls k$, $1\ls s\ls \ell$, and $t\ls \min\{r,s\}$, we have $I_t(X_{r\times s})\in \mathcal{C}_f$ and $I_t(X^{r\times s}) \in \mathcal{C}_f$.
\end{theorem}

\begin{proof}
	%We start by noticing that $I_{\omega}$ can be written as a finite sum of ideals of the form $I_t(X_{p\times q})$ where $t \leq \MIN (k,l)$, $p \leq k$ and $q \leq l$. Thus, if we prove that each of these summands is a Knutson ideal associated with $f$,  so it is $I_{\omega}$ by definition. 
	%For simplicity, we prove the result only in the case $k=\ell$, but the same proof works for the general case. 
	If $r=k$, or $s=\ell$, then $I_t(X_{r\times s}) \in \mathcal{C}_f$ \cite[Lemma 2.2]{seccia2022knutson}. Hence, we can assume   $r\neq k$ and $s \neq \ell$. We set $m:=\min\{r,s\}$ and $M:=\max\{r,s\}$.  %We can also assume without loss of generality that $r \leq s$. %(the case $q \leq p$ is symmetric). 
	We proceed by induction on $\delta:=m-t$.
	
	{\bf Base case:} Assume $\delta=0$. We define, 
	$$D_1:=\left( \deta{X_{m \times m}}, \ldots,\deta{X_{M \times M}}\right)\text{ if }M\ls k, \text{and}$$
	$$D_1:=\left( \deta{X_{m \times m}}, \ldots,\deta{X_{k \times k}}, \deta{X_{[2,k+1]}},\ldots, \deta{X_{[M-k+1,M]}} \right) \text{ otherwise}.$$
	%If $n:=\min\{k,M}$, then 
%$$I_m(X_{r \times s}) \supseteq D_1:=\left( \deta{X_{m \times m}},\deta{X_{m+1 \times m+1}}, \ldots,\deta{X_{n \times n}}, \deta{X_{[2,k+1]}},\cdots, \deta{X_{[M-k+1,M]}} \right).$$
The ideal $D_1$ is a complete intersection of  height  $M-m+1=\height\left(I_m(X_{r \times s})\right)$ and it  is contained in $I_m(X_{r \times s})$, therefore $I_m(X_{r \times s}) \in \MIN( D_1).$  
Moreover, $D_1 \in \mathcal{C}_f$ because it is generated by factors of $f$ and so it is a sum of minimal primes of $(f)$.  Thus,  $I_m(X_{r \times s})\in \mathcal{C}_f$ for every $r$ and $s$.

{\bf Induction step:}  Assume $\delta\gs 1$. We have the inclusion
\begin{equation}\label{cont1}
	I_{t}(X_{r \times s}) \supseteq D_2:=I_{t}(X_{r-1 \times s}) + I_{t}(X_{r \times s-1}) +I_{t+2}(X_{r+1 \times s+1}).
\end{equation}
The ideal  $J:= I_{t}(X_{r-1 \times s}) + I_{t}(X_{r \times s-1})$ is  ladder determinantal  (see  \autoref{sub_ladder}), so it is prime of height $(r-t+1)(s-t+1)-1$ \cite{GorlaMixed}. Moreover, $J$ does not contain $I_{t+2}(X_{r+1 \times s+1})$, then  
$$(r-t+1)(s-t+1)=\height\left(I_{t}(X_{r \times s})\right)\gs\height(D_2)\gs \height(J)+1\gs (r-t+1)(s-t+1),$$ this implies 
\begin{equation}\label{inMin}
	I_{t}(X_{r \times s}) \in \MIN ( D_2).
\end{equation}
Now, assume $s\gs r$, i.e,, $m=r$. By  the case $\delta -1$ of the induction $ I_{t}(X_{r-1 \times s})$ and $I_{t+2}(X_{r+1 \times s+1})$ belong to $\mathcal{C}_f$. 
If $r=s$, we have $I_{t}(X_{r \times s-1})\in \mathcal{C}_f$ again  by   the case $\delta -1$ of the induction, and then $I_{t}(X_{r \times s})\in \mathcal{C}_f$ by \eqref{cont1} and \eqref{inMin}. Now,  using \eqref{cont1} and \eqref{inMin} repeatedly we conclude  $I_{t}(X_{r \times s}) \in \mathcal{C}_f$ for any $s\gs r$. The case $r\gs s$ follows simmilarly by reversing the roles of $r$ and $s$.  Thus, $I_t(X_{r\times s}) \in \mathcal{C}_f$ as we wanted to show.

Likewise one shows that $I_t(X^{r\times s}) \in \mathcal{C}_f$,  finishing the proof.
\end{proof}

%To do so, we first need to specify the polynomial $f$ defining the family of Knutson ideals we are considering. 

%\begin{definition}
%\textcolor{blue}{Given a permutation $\omega \in M_{kl}$, its Schubert determinantal ideal $I_{\omega}$ can be written a sum of upper-left submatrices of the generic matrix of size $k \times l$. This suggests to choose $f$ as in the case of determinantal ideals of generic matrices (see \cite[Theorem 2.1]), as explained  below. }
%\end{definition}



As a consequence, we conclude  that Schubert determinantal ideals are Knutson ideals for the same choice of $f$, which recovers a result by Knutson \cite{knutson2009frobenius}.

\begin{corollary}\label{c:msk}
Assume \autoref{setup_Knutson_section} and  let $f:=f(X) $ be as in \autoref{def_f(X)}. Let $w \in \mathscr{M}_{k\times \ell}(\mathbb{\kk})$ be a partial  permutation and let $I_{w}$ be its corresponding Schubert determinantal ideal. Then $I_{w} \in \mathcal{C}_f$. In particular,  
$\init_\prec(I_{w})$ is a squarefree monomial ideal for any antidiagonal term order $\prec$ and 
Schubert determinantal  ideals define $F$-pure rings in positive characteristic.
\end{corollary}

\begin{proof}
Notice that $I_{w}$ can be written as a finite sum of ideals of the form $I_t(X_{r\times s})$. Thus, % where $t \leq \min(k,\ell)$, $r \leq k$ and $s \leq \ell$. 
the result follows from  \autoref{p:msk}. % each of these summands is a Knutson ideal associated with $f$,  so $I_{\omega} \in \mathcal{C}_f$ by definition.
\end{proof}

\begin{remark}
{\it Alternating sign matrix varieties}  and their defining ideals are generalizations of matrix Schubert varieties (see  \cite{klein2021bumpless} for more details).  \autoref{p:msk} applies to these ideals to show that  they belong to  $\mathcal{C}_f$. In particular, their initial ideals with respect to  antidiagonal term orders are squarefree and alternating sign matrix ideals define $F$-pure rings in positive characteristic.
\end{remark}

The next result about poset ideals generalizes \cite[Remark 79]{BSV} from maximal minors to minors of any size.

\begin{corollary}\label{c:pif}
Assume \autoref{setup_Knutson_section} and  let $f:=f(X) $ be as in \autoref{def_f(X)}. For any  ideal $\Omega\subseteq \Pi$, the poset ideal $\Omega\kk[X]$  is in $\C_{f}$. In particular, $\init_\prec(\Omega \kk[X])$ is a squarefree monomial ideal for any antidiagonal term order $\prec$ and $\kk[X]/\Omega \kk[X]$ is $F$-pure  in positive characteristic.
\end{corollary}
\begin{proof}
Given $\delta\in\Pi$ we set $\Omega_{\delta}=\{\pi\in\Pi:\pi\not\geq \delta\}\subseteq \Pi$. Notice that $\Omega_{\delta}$ is an ideal of $\Pi$, and we have the equality
$\Omega=\bigcap_{\delta\in\min\{\Pi\setminus \Omega\}}\Omega_{\delta}.$ 
The, we also have 
$\Omega \kk[X]=\bigcap_{\delta\in\min\{\Pi\setminus \Omega\}}(\Omega_{\delta}\kk[X])$ \cite[Proposition 5.2]{BookDet}. 
If $\delta=[i_1\ldots i_r|j_1\ldots j_r]$, then
\[\Omega_{\delta}\kk[X]=\sum_{t=1}^r\left(I_t\left(X_{[1,j_t-1]}\right)+I_t\left(X^{[1,i_t-1]}\right)\right)+I_{r+1}(X).\]
Therefore, the ideals $\Omega_{\delta}\kk[X]$ belong to $\C_f$ by  \autoref{p:msk}. Thus, their intersection $\Omega \kk[X]$ belongs to $\C_f$ as well. Finally, the conclusion follows as  $\Omega \kk[X]$ is the intersection of ideals in $\C_f$.
%sum of ideals in $\C_f$:
%\[\Omega \kk[X]=\sum_{\delta \in\omega}\Omega_{\delta}\kk[X].\]
\end{proof}

\begin{remark}
There are graded algebras with straightening law over fields of positive characteristic that are not $F$-pure \cite[Remark 5.2]{KoleyVarbaro}, which gives a counterexample  to a conjecture stated in page 245 of \cite{eisenbud80}. However, \autoref{c:pif} confirms the conjecture in many situations.
\end{remark}

\begin{remark}
The facts that for any
antidiagonal term order $\prec$ the ideals $\init_\prec(I_w)$ \cite[Theorem B]{knutson2005grobner}
%for any partial perrmulation $w\in\mathscr{M}_{k\times \ell}(\kk)$   
and  $\init_\prec(\Omega \kk[X])$ \cite[Theorem 2.5]{HerzogTrung} %for any poset ideal $\Omega\subset \Pi$ 
are squarefree monomial
were already known. 
\end{remark}


We introduce a new notion generalizing that of poset ideal. Given $1\leq r\leq k, 1\leq s\leq l$, we set $\Pi_{r,s}$ to be the set of all minors of $X_{r\times s}$.
\begin{definition}
A {\it generalized  ideal} of $\Pi$ is a subset $\Omega\subseteq \Pi$ such that 
$$ \text{for all }[i_1\cdots i_r|j_1\cdots j_r]\in\Omega \text{ and }\pi\in\Pi_{i_r,j_r}, \text{ if } \pi\ls [i_1\cdots i_r|j_1\cdots j_r] \text{ then } \pi\in\Omega.$$
\end{definition}

Note that  any ideal of $\Pi$ is a generalized ideal, but there exist  generalized  ideals which are not  ideals themselves, like the sets $\Pi_{r,s}$. 

Since a generalized poset ideal is the sum of  ideals of  form $\Pi_{r,s}$ (where $r,s$ vary),  the proof of \autoref{c:pif} can be adapted to show the  %and  \autoref{p:msk} 
%yields the
following result.

\begin{corollary}\label{c:gpif}
Assume \autoref{setup_Knutson_section} and  let $f:=f(X) $ be as in \autoref{def_f(X)}. For any generalized ideal $\Omega\subseteq \Pi$, the generalized poset ideal $\Omega \kk[X]$  is in $\C_{f}$. 
In particular, $\init_\prec(\Omega \kk[X])$ is a squarefree monomial ideal for any antidiagonal term order $\prec$, and $\kk[X]/\Omega \kk[X]$ is $F$-pure  in positive characteristic.
%In particular, $\init(\Omega K[X])$ is a squarefree monomial ideal for all anti-diagonal monomial order, and $K[X]/\Omega K[X]$ is $F$-pure whenever $K$ has positive characteristic.
\end{corollary}

%\section*{Acknowledgments}
%The first author was partially supported by the PRIN 2020 project 2020355B8Y ``Squarefree Gröbner degenerations, special varieties and related topics''. 
%The second author was partially funded by NSF Grant DMS \#2001645/2303605. 
% The third author is  supported by CONACyT Grant \#284598.


\bibliographystyle{alpha}
\bibliography{References}






\end{document}


\def\cprime{$'$} \def\cprime{$'$}
% \bib, bibdiv, biblist are defined by the amsrefs package.
\begin{bibdiv}
\begin{biblist}

\bib{abhyankar1988enumerative}{book}{
      author={Abhyankar, Shreeram~Shankar},
       title={Enumerative combinatorics of young tableaux},
   publisher={Marcel Dekker, Inc.},
        date={1988},
}

\bib{BSV}{article}{
      author={Benedetti, Bruno},
      author={Seccia, Lisa},
      author={Varbaro, Matteo},
       title={Hamiltonian paths, unit-interval complexes, and determinantal
  facet ideals},
        date={2022},
        ISSN={0196-8858},
     journal={Adv. in Appl. Math.},
      volume={141},
       pages={Paper No. 102407, 55},
         url={https://doi.org/10.1016/j.aam.2022.102407},
      review={\MR{4456796}},
}

\bib{Brion_Kumar_Book}{book}{
      author={Brion, Michel},
      author={Kumar, Shrawan},
       title={Frobenius splitting methods in geometry and representation
  theory},
      series={Progress in Mathematics},
   publisher={Birkh\"{a}user Boston, Inc., Boston, MA},
        date={2005},
      volume={231},
        ISBN={0-8176-4191-2},
      review={\MR{2107324}},
}

\bib{brion2007frobenius}{book}{
      author={Brion, Michel},
      author={Kumar, Shrawan},
       title={Frobenius splitting methods in geometry and representation
  theory},
   publisher={Springer Science \& Business Media},
        date={2007},
      volume={231},
}

\bib{brunsconca98}{article}{
      author={Bruns, Winfried},
      author={Conca, Aldo},
       title={K{RS} and powers of determinantal ideals},
        date={1998},
        ISSN={0010-437X},
     journal={Compositio Math.},
      volume={111},
      number={1},
       pages={111\ndash 122},
         url={https://doi.org/10.1023/A:1000287107308},
      review={\MR{1611055}},
}

\bib{bruns2022determinants}{book}{
      author={Bruns, Winfried},
      author={Conca, Aldo},
      author={Raicu, Claudiu},
      author={Varbaro, Matteo},
       title={Determinants, gr{\"o}bner bases and cohomology},
   publisher={Springer},
        date={2022},
}

\bib{BookDet}{book}{
      author={Bruns, Winfried},
      author={Vetter, Udo},
       title={Determinantal rings},
      series={Lecture Notes in Mathematics},
   publisher={Springer-Verlag, Berlin},
        date={1988},
      volume={1327},
        ISBN={3-540-19468-1},
         url={https://doi.org/10.1007/BFb0080378},
      review={\MR{953963}},
}

\bib{conca1995ladder}{article}{
      author={Conca, Aldo},
       title={Ladder determinantal rings},
        date={1995},
     journal={Journal of Pure and Applied Algebra},
      volume={98},
      number={2},
       pages={119\ndash 134},
}

\bib{concaherzog1997}{article}{
      author={Conca, Aldo},
      author={Herzog, J\"{u}rgen},
       title={Ladder determinantal rings have rational singularities},
        date={1997},
        ISSN={0001-8708},
     journal={Adv. Math.},
      volume={132},
      number={1},
       pages={120\ndash 147},
         url={https://doi.org/10.1006/aima.1997.1679},
      review={\MR{1488240}},
}

\bib{conca2020square}{article}{
      author={Conca, Aldo},
      author={Varbaro, Matteo},
       title={Square-free gr{\"o}bner degenerations},
        date={2020},
     journal={Inventiones mathematicae},
      volume={221},
      number={3},
       pages={713\ndash 730},
}

\bib{DSMNB}{article}{
      author={De~Stefani, Alessandro},
      author={Monta{\~n}o, Jonathan},
      author={N{\'u}{\~n}ez-Betancourt, Luis},
       title={Blowup algebras of determinantal ideals in prime characteristic},
        date={2021},
     journal={arXiv preprint arXiv:2109.00592},
}

\bib{eisenbud80}{inproceedings}{
      author={Eisenbud, David},
       title={Introduction to algebras with straightening laws},
        date={1980},
   booktitle={Ring theory and algebra, {III} ({P}roc. {T}hird {C}onf., {U}niv.
  {O}klahoma, {N}orman, {O}kla., 1979)},
      series={Lecture Notes in Pure and Appl. Math.},
      volume={55},
   publisher={Dekker, New York},
       pages={243\ndash 268},
      review={\MR{584614}},
}

\bib{fulton1992flags}{article}{
      author={Fulton, William},
       title={Flags, schubert polynomials, degeneracy loci, and determinantal
  formulas},
        date={1992},
     journal={Duke Math. J.},
      volume={68},
      number={1},
       pages={381\ndash 420},
}

\bib{gao2022minimal}{article}{
      author={Gao, Shiliang},
      author={Yong, Alexander},
       title={Minimal equations for matrix schubert varieties},
        date={2022},
     journal={arXiv preprint arXiv:2201.06522},
}

\bib{GS}{article}{
      author={Glassbrenner, Donna},
      author={Smith, Karen~E.},
       title={Singularities of certain ladder determinantal varieties},
        date={1995},
        ISSN={0022-4049},
     journal={J. Pure Appl. Algebra},
      volume={101},
      number={1},
       pages={59\ndash 75},
         url={http://dx.doi.org/10.1016/0022-4049(94)00038-K},
      review={\MR{1346428 (96h:13029)}},
}

\bib{gonciulea2000mixed}{article}{
      author={Gonciulea, Nicolae},
      author={Miller, Claudia},
       title={Mixed ladder determinantal varieties},
        date={2000},
     journal={Journal of Algebra},
      volume={231},
      number={1},
       pages={104\ndash 137},
}

\bib{GorlaMixed}{article}{
      author={Gorla, Elisa},
       title={Mixed ladder determinantal varieties from two-sided ladders},
        date={2007},
        ISSN={0022-4049},
     journal={J. Pure Appl. Algebra},
      volume={211},
      number={2},
       pages={433\ndash 444},
         url={https://doi.org/10.1016/j.jpaa.2007.01.016},
      review={\MR{2340461}},
}

\bib{GrifoHuneke}{article}{
      author={Grifo, Elo\'{\i}sa},
      author={Huneke, Craig},
       title={Symbolic powers of ideals defining {F}-pure and strongly
  {F}-regular rings},
        date={2019},
        ISSN={1073-7928},
     journal={Int. Math. Res. Not. IMRN},
      number={10},
       pages={2999\ndash 3014},
         url={https://doi.org/10.1093/imrn/rnx213},
      review={\MR{3952556}},
}

\bib{BinomialSumSP}{article}{
      author={H{\`a}, Huy~T{\`ai}},
      author={Nguyen, Hop~Dang},
      author={Trung, Ngo~Viet},
      author={Trung, Tran~Nam},
       title={Symbolic powers of sums of ideals},
        date={2020},
        ISSN={0025-5874},
     journal={Math. Z.},
      volume={294},
      number={3-4},
       pages={1499\ndash 1520},
  url={https://doi-org.biblio.cimat.remotexs.co/10.1007/s00209-019-02323-8},
      review={\MR{4074049}},
}

\bib{H98}{article}{
      author={Hara, Nobuo},
       title={A characterization of rational singularities in terms of
  injectivity of {F}robenius maps},
        date={1998},
        ISSN={0002-9327},
     journal={Amer. J. Math.},
      volume={120},
      number={5},
       pages={981\ndash 996},
  url={http://muse.jhu.edu/journals/american_journal_of_mathematics/v120/120.5hara.pdf},
      review={\MR{1646049}},
}

\bib{HerzogTrung}{article}{
      author={Herzog, J\"{u}rgen},
      author={Trung, Ng\^{o}~Vi\^{e}t},
       title={Gr\"{o}bner bases and multiplicity of determinantal and
  {P}faffian ideals},
        date={1992},
        ISSN={0001-8708},
     journal={Adv. Math.},
      volume={96},
      number={1},
       pages={1\ndash 37},
         url={https://doi.org/10.1016/0001-8708(92)90050-U},
      review={\MR{1185786}},
}

\bib{HHStrongFreg}{article}{
      author={Hochster, Melvin},
      author={Huneke, Craig},
       title={Tight closure and strong {$F$}-regularity},
        date={1989},
        ISSN={0037-9484},
     journal={M\'em. Soc. Math. France (N.S.)},
      number={38},
       pages={119\ndash 133},
        note={Colloque en l'honneur de Pierre Samuel (Orsay, 1987)},
      review={\MR{1044348}},
}

\bib{klein2021bumpless}{article}{
      author={Klein, Patricia},
      author={Weigandt, Anna},
       title={Bumpless pipe dreams encode {G}r\"obner geometry of {S}chubert
  polynomials},
        date={2021},
     journal={arXiv preprint arXiv:2108.08370},
}

\bib{knutson2009frobenius}{article}{
      author={Knutson, Allen},
       title={Frobenius splitting, point-counting, and degeneration},
        date={2009},
     journal={arXiv preprint arXiv:0911.4941},
}

\bib{knutson2005grobner}{article}{
      author={Knutson, Allen},
      author={Miller, Ezra},
       title={Gr{\"o}bner geometry of schubert polynomials},
        date={2005},
     journal={Annals of Mathematics},
       pages={1245\ndash 1318},
}

\bib{KoleyVarbaro}{article}{
      author={Koley, Mitra},
      author={Varbaro, Matteo},
       title={Gr{\"o}bner deformation and {$F$}-singularities},
        date={2021},
     journal={arXiv preprint arXiv:2107.12116, to appear in Math. Nachr.},
}

\bib{ma2021f}{article}{
      author={Ma, Linquan},
      author={Polstra, Thomas},
       title={F-singularities: a commutative algebra approach},
        date={2021},
     journal={Preprint, available at https://www. math. purdue. edu/\~{}
  ma326/F-singularitiesBook. pdf},
}

\bib{mehta1985frobenius}{article}{
      author={Mehta, Vikram~B},
      author={Ramanathan, Annamalai},
       title={Frobenius splitting and cohomology vanishing for schubert
  varieties},
        date={1985},
     journal={Annals of Mathematics},
      volume={122},
      number={1},
       pages={27\ndash 40},
}

\bib{MillerSturmfels}{book}{
      author={Miller, Ezra},
      author={Sturmfels, Bernd},
       title={Combinatorial commutative algebra},
      series={Graduate Texts in Mathematics},
   publisher={Springer-Verlag},
     address={New York},
        date={2005},
      volume={227},
        ISBN={0-387-22356-8},
      review={\MR{2110098 (2006d:13001)}},
}

\bib{procesi2007lie}{book}{
      author={Procesi, Claudio},
       title={Lie groups: an approach through invariants and representations},
   publisher={Springer},
        date={2007},
      volume={115},
}

\bib{seccia2021knutson}{article}{
      author={Seccia, Lisa},
       title={Knutson ideals and determinantal ideals of hankel matrices},
        date={2021},
     journal={Journal of Pure and Applied Algebra},
      volume={225},
      number={12},
       pages={106788},
}

\bib{seccia2022knutson}{article}{
      author={Seccia, Lisa},
       title={Knutson ideals of generic matrices},
        date={2022},
     journal={Proceedings of the American Mathematical Society},
      volume={150},
      number={05},
       pages={1967\ndash 1973},
}

\bib{SmithFrational}{article}{
      author={Smith, Karen~E.},
       title={{$F$}-rational rings have rational singularities},
        date={1997},
        ISSN={0002-9327},
     journal={Amer. J. Math.},
      volume={119},
      number={1},
       pages={159\ndash 180},
  url={http://muse.jhu.edu.proxy.lib.umich.edu/journals/american_journal_of_mathematics/v119/119.1smith.pdf},
      review={\MR{1428062 (97k:13004)}},
}

\bib{Sull}{article}{
      author={Sullivant, Seth},
       title={Combinatorial symbolic powers},
        date={2008},
        ISSN={0021-8693},
     journal={J. Algebra},
      volume={319},
      number={1},
       pages={115\ndash 142},
         url={https://doi.org/10.1016/j.jalgebra.2007.09.024},
      review={\MR{2378064}},
}

\end{biblist}
\end{bibdiv}
